\section{Source, scope, and conventions}

The supplied final LaTeX and PDF are the manuscript reference for this
release. Statement and equation numbers refer to that version.
See the project statement map for declaration-level correspondence and
the assumption catalogue for quoted theorems and bridges.

The displayed statements and their proof routes describe the mathematics
of the paper. A checked Lean node names its actual declaration and
identifies the precise part of the argument it proves. An open node
records a proof obligation. Quoted inputs belong to the assumption
catalogue, and the estimates newly proved in the manuscript remain
obligations of this formalization.

\printgraph

\section{Finite counting and the selector}

\begin{definition}[Configuration and arithmetic rescaling]\label{d:Cq}
Write $[N]=\{1,\ldots,N\}$ and $P_0=0$. The original integer
polynomials satisfy $P_i(0)=0$ and
$1\le \deg P_1<\cdots<\deg P_\ell=d$. A configuration is an $x\in\Z$
and positive integer $n$ such that every $x+P_i(n)$ belongs to $A$,
and the map from $i\in\{0,\ldots,\ell\}$ to these points is injective.
For a positive integer $q$ put $P_{i,q}(y)=P_i(qy)/q$.
For $P(y)=a_ey^e+\cdots+a_1y$, $e\ge1$, the $(C,q)$-coefficient
condition is
\[
 |a_j|\le C|a_e|\quad(1\le j<e),\qquad
 a_e=a'_e q^{e-1},\qquad a'_e\in\Z,\quad 0<|a'_e|\le C.
\]
For $R_0=0$ and a given family $R_1,\ldots,R_{k-1}$, define
\[
 T^{\mathbf R}_{N,M,v}(\mathbf f)
 =\frac{1}{NM}\sum_{x\in\Z}\sum_{y\in[M]}
       v(y)\prod_{i=0}^{k-1} f_i(x+R_i(y)).
\]
\end{definition}
\source{d:Cq}

The normalization uses the spatial scale $N$ and the parameter length $M$.
The support of $f_0$ bounds the spatial sum because $R_0=0$.
From $P_i(0)=0$, the coefficient of $y^j$ in $P_{i,q}$ is $q^{j-1}$
times that in $P_i$. This gives the coefficient condition with a fixed
$C$ for every $q\ge1$. The same arithmetic gives the successive
rescalings. If $q'=q^{B+1}$ then
\[
 q^B P_{i,q'}(y)=P_{i,q}(q^B y),\qquad
 q'y+1=q(q^B y)+1.
\]
An affine embedding $x\mapsto a+q^B x$ therefore preserves the parameter
prime, its selector value, and pairwise distinctness.

\begin{definition}[Phase tests and Gowers normalization]\label{phase-tests}
For a finite nonempty progression $P$ and real polynomial tests, put
\[
 \norm{v}_{u^s[P]}
  =\sup_{\substack{\phi\in\R[y]\\\deg\phi<s}}
    \abs{\frac1{|P|}\sum_{y\in P}v(y)\ee(-\phi(y))},
 \qquad \ee(t)=e^{2\pi it}.
\]
For a finitely supported $g:\Z\to\C$, let
\[
 \norm{g}^{2^s}_{U^s(\Z)}
  =\sum_{x,h_1,\ldots,h_s\in\Z}
     \prod_{\omega\in\{0,1\}^s}
       \mathcal C^{|\omega|}g(x+\omega\cdot h),
 \qquad
 \norm{g}_{U^s[I]}
 =\frac{\norm{g\one_I}_{U^s(\Z)}}{\norm{\one_I}_{U^s(\Z)}}.
\]
Here $\mathcal C$ is complex conjugation and $I$ is a finite
nonempty interval or progression. For $s=1$ this is the absolute
average. Normalized monotonicity bounds $U^r[I]$ by $U^s[I]$
up to constants depending only on $r,s$, for $1\le r\le s$.
Multiplication by $\ee(\phi)$ with $\deg\phi<s$ preserves $U^s[I]$,
so these facts also give $u^s[I]\ll_s U^s[I]$.
\end{definition}

\paragraph{Proof of the finite norm estimates.}
For each cube coordinate $i$, reflect the array of vertex functions
onto the two $i$-faces. The exact substitution $y=x+h_i$ factors the
sum into a correlation of two functions of the other directions. The
zero extension allows both base sums to run over the same interval.
Cauchy--Schwarz bounds the squared mixed form by the two reflected
forms. Repeating in all coordinates bounds the mixed form by the
product of the vertex norms. Expanding the vertices for $f+g$
gives the triangle inequality. The real cube power $Q_s(f)$ is also
a sum of squared magnitudes of the $(s-1)$-direction correlations.
Viewing $Q_s(f)$ as the sum of the $s$-direction correlations and
using Cauchy--Schwarz gives $Q_s(f)^2\le(2W+1)^sQ_{s+1}(f)$ for
an interval of width $W$. The indicator cube count lies between
$(W/(s+1))^{s+1}$ and $W(2W)^s$. For the lower bound take the base
point and $s$ nonnegative directions in intervals of length
$\lfloor(W-1)/(s+1)\rfloor+1$. The base and the singleton vertices
bound the upper count. Division by the indicator powers therefore
gives a uniform adjacent-order bound.

An affine bijection from $[M]$ to a positive-step progression matches
every surviving cube because each direction is divisible by the step.
The same bijection transports the indicator and takes ordinary
polynomials of degree below $s$ to polynomials of degree below $s$.
The norm at order one is the magnitude of the mean. Phase invariance
and order comparison then bound each phase test before the supremum.
For dilation on $[b,L+b]$, use the $q$ linear characters of
$\mathbb Z/q\mathbb Z$ to restrict to the residue class $b$. Their
coefficients have total absolute mass one. Remove the ambient point
$b$ using the triangle inequality and bound its singleton norm using
the all-zero directions. Exact affine cube transport and the volume
bounds compare the ambient width $L+1$ with $\lfloor L/q\rfloor$.
Their ratio is at most $2q$, which gives the power $q^{(s+1)/2^s}$.
All steps are checked in Lemma~\ref{checked-norm-foundations}.

\paragraph{Circle Fourier summability.}
For a $K$-Lipschitz cutoff $\chi$ with coefficients $c_h$, the difference
$\chi(\cdot+t)-\chi$ has $L^2$ norm at most $K|t|$ in unit Haar
measure. Its $h$-th coefficient is $(\ee(ht)-1)c_h$.
Parseval bounds the sum of squared moduli on each finite frequency
set. Divide by $t^2$ for $t\ne0$ and let $t\to0$ on that finite set.
This gives a uniform bound for the partial sums of
$\sum_h (2\pi h)^2|c_h|^2$. Cauchy--Schwarz against the summable
squares $h^{-2}$ bounds $\sum_{h\ne0}|c_h|$. The zero coefficient
is at most $\norm{\chi}_\infty$. Mathlib's Fourier reconstruction
then identifies the absolutely convergent series with $\chi$.

\begin{definition}[Heisenberg and bracket selectors]\label{selector}
Put
\[
 \theta=\frac{\alpha\beta}{2},\qquad
 Z(n)=\theta n(n-1)-\alpha n\lfloor\beta n\rfloor,
 \qquad Z_{\rm br}(n)=\alpha(n+1)\lfloor\beta(n+1)\rfloor .
\]
The assumptions are $c_\alpha=\inf_{m\ge1}m\tn{m\alpha}>0$ and
$\beta\in(0,1)$ with an increasing sequence of reduced rationals
$a_j/r_j$ of increasing prime denominator $r_j$, tending to $\beta$,
such that
\[
 r_{j+1}\asymp\exp(\exp(\sqrt{\log r_j})),\qquad
 0<\beta-a_j/r_j\ll r_{j+1}^{-1}.
\]
Write $\T=\R/\Z$ and let $J$ be an open arc of length in $(0,1)$.
The return set is
$\mathcal B_{Z,J}=\{n\in\Z_{>0}:Z(n)\bmod1\in J\}$.
For $V=Z$ or $Z_{\rm br}$ set
\[
 \mathcal P_{V,J}=\{p\in\PP:V(p-1)\bmod1\in J\},\qquad
 \pi_{V,J}(N)=|\mathcal P_{V,J}\cap[N]|.
\]
Set $\Lambda_q(y)=(\varphi(q)/q)\Lambda(qy+1)$,
$F_q(y)=F(Z(qy))$ and $\mu_F=\int_\T F$.
\end{definition}

Let $H=\R^3$ have multiplication
$(x,y,z)(x',y',z')=(x+x',y+y',z+z'+xy')$.
With $\Gamma=\Z^3$ and $g_0=(\alpha,\beta,0)$, direct multiplication
gives $g_0^n=(n\alpha,n\beta,\theta n(n-1))$ for all integer $n$.
Right multiplication by $(u,v,w)\in\Gamma$ changes
$z-x\lfloor y\rfloor$ by $w-u(\lfloor y\rfloor+v)\in\Z$.
Thus on the chart $0<\{y\}<1$, the vertical phase
$z-x\lfloor y\rfloor\bmod1$ gives $Z(n)\bmod1$ on the orbit.
The Haar probability measure projects to uniform measure on the
horizontal torus, and vertical fibre measure is uniform.
For $n\ge1$, irrationality of $\beta$ ensures $\{\beta n\}>0$,
so these are returns to the open vertical window of measure $|J|$.

Admissible $\beta$ are dense in $(0,1)$. To see this, start with a
reduced rational in the middle third of a prescribed interval, with
large prime denominator $r_1$. Bertrand's postulate supplies a
prime $r_{j+1}$ comparable to the prescribed growth function.
Choose $a_{j+1}=\lfloor r_{j+1}a_j/r_j\rfloor+1$.
The successive positive increments are less than $1/r_{j+1}$.
The denominator growth bounds the entire tail by a constant times
its first term. A large $r_1$ keeps the fractions within the prescribed
interval. This also gives $\beta-a_j/r_j=o(r_j^{-D})$ for every
fixed $D$. The positive errors and rational spacing exclude rational
limits, giving the Liouville property. Lean checks this density construction
and the implication from admissibility to the Liouville property.

Proposition~\ref{checked-intro-fast-obstruction} gives the introductory
counterexample using faster prime-denominator approximations.

\begin{proposition}[Quoted inputs]\label{external-inputs}
\lean{NilBohr.KMTTShiftedPrimeDensityBound,NilBohr.KMTTLocalInverse,NilBohr.KMTTLocalInverseDegreeUniform,NilBohr.KMTTWeightedGowers,NilBohr.KMTTAmbientPrimeModel,NilBohr.KMTTRelativeCramer,NilBohr.BanksRealZeroBound,NilBohr.GreenTaoQuantitativeLeibman,NilBohr.MertensIntervalProductBound,NilBohr.FundamentalSieveLemma,NilBohr.TeravainenFixedFamily,NilBohr.MirekWanWrightPointwise,NilBohr.TeravainenWangSparseGreenTao}
The independent analytic inputs include the KMTT prime-model estimates
(Lemmas 3.1 and 4.3(i)), its generalized von Neumann estimate (Proposition 3.3),
inverse estimate (Lemma 3.4), and
shifted-prime theorem (Theorem 1.2). The other analytic inputs are
Banks's explicit zero-free bound (Theorem 1.1), Green--Tao quantitative
Leibman (Theorem 2.9 and its erratum), the dimension-one fundamental
sieve lemma and Mertens interval bound (Koukoulopoulos, Theorems 19.1
and 3.4(c)), Ter\"av\"ainen's fixed-family
estimate (Theorem 4.1), and
Mirek--Wan--Wright pointwise convergence (Theorem 1.1(ii)).
Ter\"av\"ainen--Wang's finite density bound
(Theorem~1.1) supplies both progression implications.
The degree-uniform KMTT inverse interface fixes a bound
for the modulus exponent before the coefficient constant or family
length is chosen.
\end{proposition}

The project assumption catalogue gives the source versions, individual
interfaces, and status of the translation obligations. The scalar
and product obstructions are deduced from the quantitative Leibman
interface in Lemma~\ref{checked-green-tao-adapter}.
The Gowers, dilation and circle Fourier premises are supplied by
Lemma~\ref{checked-norm-foundations}. Lemma~\ref{checked-weak-regularity}
derives weak regularity from the inverse estimate. A single classical prime number
theorem limit suffices by Lemma~\ref{checked-prime-normalization}.
The ambient prime model and sieve premises supply that limit, its
fixed-modulus version and Vinogradov's linear equidistribution
interface by Lemma~\ref{checked-prime-distribution}.


\section{Local extraction of polynomial phases}

\begin{lemma}[5.1, inverse estimate with a common modulus]\label{p:localinverse}
\lean{NilBohr.WeightedLocalInverseEstimate,NilBohr.weighted_local_inverse_of_kmtt_degree_uniform,NilBohr.weighted_local_inverse_of_kmtt}\leanok
\uses{d:Cq}
Fix $C,d,k$, $k\ge2$, and let $N\ge3$,
$1\le q\le\exp((\log N)^{1/10})$ and
$M=\lfloor(N/q^{d-1})^{1/d}\rfloor\ge1$. Assume the shifts
have $(C,q)$-coefficients, distinct increasing degrees ending in $d$,
and the $f_i$ are $1$-bounded with support in $[-CN,CN]$.
If $|v|\le1$ and $|T_{N,M,v}(\mathbf f)|\ge\delta$ with
$0<\delta<1/2$, then either $N\ll\delta^{-O(1)}$, or there is a fixed
$L=q^B$, $B=O_d(1)$, and for each $i\ge1$ an integer $V_i$
and a unit phase $\Psi_i$ with
\[
 M^{\deg P_i}(q/\delta)^{-O(1)}\ll V_i\le M^{\deg P_i},\quad
 LV_i\le N,\qquad
 \abs{\sum_x f_i(x)\Psi_i(x)}\gg\delta^{O(1)}N.
\]
The phase is linear on each intersection of a block
$((j-1)LV_i,jLV_i]$ and a residue class modulo $L$.
\end{lemma}
\source{p:localinverse}
\begin{proof}[Proof in the manuscript]
\proves{p:localinverse}\uses{external-inputs,checked-difference}\leanok
The first step removes the weight while retaining a polynomial
fraction of spatial differences. By Cauchy--Schwarz,
\[
 \delta^2N^2\le\E_{y\in[M]}
       \abs{\sum_x\prod_{i=0}^{k-1} f_i(x+P_i(y))}^2
 =\sum_h\E_{y\in[M]}\sum_x\prod_{i=0}^{k-1}
       \Delta_h f_i(x+P_i(y)),
\]
where $\Delta_h f(x)=f(x+h)\overline{f(x)}$.
Only $O_C(N)$ terms in $h$ are nonzero and each is $O_C(N)$.
The lower bound therefore gives $\gg_C\delta^2N$ shifts $h$
with unnormalized count of size $\gg_C\delta^2N$. Translate
the spatial supports into $[N']$ with $N'\asymp_C N$ and
$N'\ge N$. The KMTT inverse theorem applies at length $M$,
and its normalized threshold is a fixed multiple of $\delta^2$.
Its small alternative $M\ll(q/\delta)^{O(1)}$ gives
$N<q^{d-1}(M+1)^d\ll(q/\delta)^{O(1)}$. As
$q\le\exp((\log N)^{1/10})$, any fixed power of $q$ is at most
$N^{1/2}$ for sufficiently large $N$, giving the required
small alternative.

In the other branch, KMTT supplies an auxiliary
$q'_h\ll\delta^{-O(1)}$, an exponent $b_h=O_d(1)$ and
$1$-bounded $\varphi_{h,i}$ such that the correlation with
$\Delta_hf_i$ is at least $\delta^{O(1)}N$ up to a fixed constant, and
\[
 |\varphi_{h,i}(x+q'_h q^{b_h} t)-\varphi_{h,i}(x)|
   \ll(q/\delta)^{O(1)}M^{-\deg P_i}|t|.
\]
Take $B$ to be a uniform bound for $b_h$ and choose $D_0\ge B$
large. Let $V_i=\lfloor M^{e_i}(q/\delta)^{-D_0}\rfloor$,
$e_i=\deg P_i$. A value inside the floor below $2$ forces the
small alternative by the same bound on powers of $q$.
Since $L(q/\delta)^{-D_0}\le1$, the inequality $LV_i\le N$ follows.
The family of full blocks meeting the support contributes
$O_C(N/V_i)$ progressions $P=a+L[V_i]$, including cells at the ends.

On each such $P$, hold $\varphi_{h,i}(a+Lt)$ constant on each
class of $t$ modulo $q'_h$ by choosing a representative in $[V_i]$
for each class that occurs. The approximation error is
$O((q/\delta)^{-D_0+O(1)})$ by replacing $t$ in the Lipschitz
bound by $q^{B-b_h}(t-t')/q'_h$. The correlation over the
support changes by at most $O(N(q/\delta)^{-D_0+O(1)})$.
Fourier inversion modulo $q'_h$ with coefficients of modulus at
most one then gives, for one $\omega_h=m/q'_h$,
\[
 \sum_{P=a+L[V_i]}
      \abs{\sum_{t\in[V_i]}\Delta_h f_i(a+Lt)\ee(\omega_h t)}
      \gg\delta^{O(1)}N .
\]
There are $\delta^{-O(1)}$ possible frequency values. For each
fixed $i$, pigeonhole one frequency $\omega$ across
$\gg\delta^{O(1)}N$ choices of $h$.

To undo differencing define
$S_P=\max_{\theta\in\T}|\sum_{t\in[V_i]}f_i(a+Lt)\ee(\theta t)|$.
If $A_P(h)=\sum_t\Delta_h f_i(a+Lt)\ee(\omega t)$, then
\[
 \widehat{A_P}(\xi)
 =\widehat f_i(\xi)\sum_t\overline{f_i(a+Lt)}
         \ee(\omega t+\xi(a+Lt)).
\]
Plancherel gives $\sum_h|A_P(h)|^2\le S_P^2\norm{f_i}_2^2$.
Cauchy--Schwarz over $O_C(N)$ shifts and $O_C(N/V_i)$ progressions
then gives $\sum_P S_P^2\gg\delta^{O(1)}NV_i$. As $S_P\le V_i$,
$\sum_P S_P\gg\delta^{O(1)}N$. On each cell choose a maximizing
frequency (compactness of $\T$) and rotate its correlation to be the
nonnegative number $S_P$. Define the phase to be $1$ on other cells.
This finishes the assertion with common modulus $L$.
\end{proof}
\noindent\textbf{Lean status.} The weighted deduction is checked
from the degree-uniform KMTT inverse interface. Lean chooses one bound
$B$ after fixing $d$ and before fixing $C$ and the number of shifts.
The argument carries this same bound through spatial differencing,
finite Fourier localization, and Plancherel. The other constants
may depend on $C,d,k$, as allowed in the manuscript.

\begin{theorem}[4.2, local generalized von Neumann]\label{t:main}
\lean{NilBohr.LocalGeneralizedVonNeumann,NilBohr.localGeneralizedVonNeumann_of_kmtt,NilBohr.localGeneralizedVonNeumann_twoInput,NilBohr.countingForm_twoInput_le}\leanok
\uses{d:Cq,phase-tests}
Under the scale, support and coefficient assumptions of Lemma~5.1,
there exist $\eta,A_1,A_2>0$ and $B\in\N$, depending only on
$C,d,k$, such that, for every $0<\rho\le1/2$,
\[
 \abs{T_{N,M,v}(\mathbf f)}
 \ll_{C,d,k}
 N^{-\eta}+\rho+
 \left(\sup_{I,b,u}\norm{v}_{u^{d+1}[I\cap(u+q^b\Z)]}\right)^\eta .
\]
The supremum ranges over integer intervals $I\subset[M]$ with
$|I|\ge Mq^{-A_1}\rho^{A_2}$ and integers $0\le b\le B$ and $u$ such
that $I\cap(u+q^b\Z)$ is nonempty. For $k=2$, use $b=0$ and $I=[M]$.
\end{theorem}
\source{t:main}
\begin{proof}[Proof in the manuscript]
\proves{t:main}\uses{p:localinverse,checked-local-phase}\leanok
For $k=2$ Fourier inversion in $x$ writes the count as
\[
 N^{-1}\int_\T\widehat f_0(-\theta)\widehat f_1(\theta)
       \E_{y\in[M]}v(y)\ee(\theta P_1(y))\,d\theta.
\]
Cauchy--Schwarz and Plancherel bound this by
$N^{-1}\norm{f_0}_2\norm{f_1}_2\norm{v}_{u^{d+1}[M]}$.

For $k\ge3$, suppose the count is at least $\delta$ in modulus.
For input $i\ge1$, with the current other inputs, write
\[
 D_i(z)=\E_{y\in[M]}v(y)\prod_{j\ne i}
          f_j(z+P_j(y)-P_i(y)),\qquad
 T=N^{-1}\sum_z f_i(z)D_i(z).
\]
The unchanged $f_0$ forces $D_i$ to be supported in a fixed
enlargement of $[-CN,CN]$, since each shift on $[M]$ is
$O_{C,d}(N)$. The modified count for $f_i=\overline D_i$ equals
$N^{-1}\sum_z|D_i(z)|^2$, which by Cauchy--Schwarz is at least a fixed
power of the previous count up to constants. Apply Lemma~5.1 to this
tuple and use $\overline\Psi_i$ as the replacement. The conjugate of
the inverse-theorem correlation is exactly the replacement count
times $N$. Truncate inserted phases to the fixed enlarged support
when applying the lemma again. The truncation changes no term with
$f_0(x)\ne0$. Repeat for all $k-1$ shifted inputs. The bound for the
common exponent depends only on $d$.

Now let $H=\lfloor M(q/\delta)^{-D_1}\rfloor$ for a fixed large
$D_1$. A value inside the floor less than $2$ forces the small
alternative as before. Otherwise divide $[M]$ into length-$H$
intervals and one tail of length less than $H$. On a complete
interval $I$ with first point $y_0$,
\[
 |P_i(y)-P_i(y_0)|\ll_{C,d}q^{e_i-1}M^{e_i-1}H .
\]
For each $i$, delete the $x$ at which $x+P_i(y_0)$ is at most this
far from a block boundary in $LV_i\Z$. There are
$O_C(N/(LV_i))$ relevant boundaries. The normalized total loss
including the tail is at most a fixed constant times
\[
 \frac HM+\sum_i\frac{q^{e_i-1}M^{e_i-1}H}{LV_i}
 \ll(q/\delta)^{-D_1+O(1)}.
\]
The retained contribution still has modulus $\gg\delta^{O(1)}$
for sufficiently large $D_1$. The weights of the retained averages
sum to at most $N^{-1}\sum_x|f_0(x)|\ll_C1$. Some retained
pair $(x,I)$ therefore gives an average of comparable modulus.

Restricted to $y\equiv u\pmod L$, every argument $x+P_i(y)$
stays in a single fixed block and residue class. The product of local
linear phases is $\ee(\phi(y))$ with
$\phi(y)=\sum_i[\theta_i(x+P_i(y))+\gamma_i]$ of degree at most $d$.
Taking absolute values in the convex combination of class averages
finds a nonempty class with correlation $\gg\delta^{D}$ and
$|I|\gg Mq^{-D}\delta^D$ for some fixed $D$.
Set $\delta$ to a small fixed multiple of the count.
If $\delta\le\rho$, use the $\rho$ term.
If $N\ll\delta^{-D}$, use the $N^{-\eta}$ term.
Otherwise choose $A_1,A_2$ to absorb the constants in the interval
length bound using $\rho\le1/2$, and take $\eta\le1/D$.
\end{proof}
\noindent\textbf{Lean status.} The theorem is checked from the
KMTT inverse estimate. The two-input proof retains the explicit support
factor $|X_0|/N$ for spatial window $X_0$. In the multi-input proof,
Lean partitions $[M]$ into $\lfloor M/H\rfloor$ full
intervals and a tail of length $M\bmod H$. If $R_i$ bounds the
displacement across one interval, the error is bounded by
$|X_0|H/(NM)+\sum_i 5(2C+4)R_i/(LV_i)$. For the first term,
the tail uses $H$ as a strict upper bound. For each other term,
the counted points are within displacement $R_i$ of a block
boundary, rounded with an integer ceiling where necessary.
Lean chooses $H$ to absorb both costs into the lower bound for
the phase-replaced count. If that choice yields $H=0$, the same
scale estimates place $N$ in the small-scale alternative. The final
conversion keeps the interval-length condition and the bound on $b$
uniform in $N,q,\rho$ and the test weights.

\section{Heisenberg cancellation and the localized multiplier}

\begin{lemma}[6.1, horizontal Diophantine separation]\label{l:heisenberg-horizontal}
\lean{NilBohr.heisenberg_horizontal,NilBohr.horizontal_scale_bounds,NilBohr.horizontal_diophantine_bound_uniform_in_beta,NilBohr.horizontal_diophantine_bound}
\leanok
\uses{selector}
Choose $r_j\le X^{1/100}<r_{j+1}$ for sufficiently large $X$.
For each fixed $L>0$, $r_j/(\log X)^L\to\infty$.
Uniformly when $1\le v\le X^{1/5000}$, $r_j\nmid v$ and
$0<\max(|k|,|l|)\le(\log X)^L$ for integers $v,k,l$, one has
$\tn{v(k\alpha+l\beta)}\gg_\alpha X^{-1/20}$.
\end{lemma}
\source{l:heisenberg-horizontal}
\begin{proof}[Proof in the manuscript]
\proves{l:heisenberg-horizontal}\uses{checked-floor}\leanok
The growth rule gives $\log\log X\le\sqrt{\log r_j}+O(1)$,
so $r_j$ dominates every fixed power of $\log X$.
Also $r_{j+1}/r_j^2\gg X^{1/200}$.
If $r_j\le X^{1/400}$ use $r_{j+1}>X^{1/100}$.
Otherwise the growth rule gives $r_{j+1}\gg r_j^4$ and suffices.

For $k\ne0$, multiplication by $r_j$ and the reverse triangle
inequality give
\[
 \tn{v(k\alpha+l\beta)}
 \ge r_j^{-1}\tn{r_jvk\alpha}
       -|l|v|\beta-a_j/r_j|
 \ge\frac{c_\alpha}{|k|v r_j^2}-O(|l|v/r_{j+1}).
\]
The error-to-main ratio is at most a constant times
$X^{1/2500-1/200}(\log X)^{2L}$, which tends to zero. The main term
dominates $X^{-1/20}$ since $r_j\le X^{1/100}$.
For $k=0$, the primality of $r_j$ and $r_j\nmid vla_j$ give
$\tn{vla_j/r_j}\ge1/r_j$. The error
$O(|l|v/r_{j+1})$ is $o(1/r_j)$. This gives the same bound.
\end{proof}
\noindent\textbf{Lean status.} Checked, including eventual selection
of the bracketing denominator and the uniformity of both estimates.
The positive lower-bound constant is chosen from $\alpha$ alone,
before the choice of $\beta$ and its approximation constants.

\begin{proposition}[4.3, uniform Heisenberg phase bounds]\label{p:heisenberg-phase}
\lean{NilBohr.heisenbergPhaseEstimates_of_quantitativeLeibman,NilBohr.scalar_heisenberg_log_saving,NilBohr.twisted_heisenberg_log_saving,NilBohr.centeredHeisenbergPhaseEstimate_of_obstructions}\leanok
\uses{selector,phase-tests}
At the scale of Lemma~6.1 let $a\in\Z$,
$1\le v\le X^{1/10000}$, $r_j\nmid v$, and $K\ge X^{1/2}$.
For fixed $A,L>0$, a positive integer $d$, and a Lipschitz function $\Phi$ on $H/\Gamma$
with $\norm{\Phi}_{\mathrm{Lip}}\le(\log X)^L$, one has
\[
 \abs{\E_{t\in[K]}\Phi(g_0^{a+vt}\Gamma)-\int\Phi}
 \ll_{A,L,\alpha,\beta}(\log X)^{-A}.
\]
With zero integral over every vertical fibre, one also has
\[
 \sup_{\deg\phi\le d}
 \abs{\E_{t\in[K]}\Phi(g_0^{a+vt}\Gamma)\ee(-\phi(t))}
 \ll_{A,L,d,\alpha,\beta}(\log X)^{-A}.
\]
For fixed Lipschitz $F$, the same twisted bound holds with
$\Phi(g_0^{a+vt}\Gamma)$ replaced by $F(Z(a+vt))-\mu_F$.
The twisted and centered statements also cover constant tests,
by restricting the degree-one range.
\end{proposition}
\source{p:heisenberg-phase}
\begin{proof}[Proof in the manuscript]
\proves{p:heisenberg-phase}
\uses{l:heisenberg-horizontal,checked-scalar,checked-twisted,external-inputs,checked-centered,checked-green-tao-adapter}\leanok
First prove ordinary equidistribution in a form that allows
dilations. For any positive integer $s$ bounded by a fixed power of
$\log X$, put $D_s(x,y,z)=(sx,sy,s^2z)$. Since
$D_s\Gamma\subset\Gamma$, $D_s$ induces a quotient map. The inverse
$D_{1/s}$ is applied to group elements only. Consider
$D_{1/s}(g_0^{a'+v't})\Gamma$ for $1\le t\le K_0$, where
$K_0\ge X^{1/3}$, $1\le v'\le X^{1/5000}$, $r_j\nmid v'$,
and $a'\in\Z$. If equidistribution with parameter
$(\log X)^{-B}$ fails, quantitative Leibman gives a nonzero
character $(k,l)$ with $|k|,|l|\le(\log X)^{O_B(1)}$ and
\[
 K_0\tn{v'(k\alpha+l\beta)/s}\ll_B(\log X)^{O_B(1)} .
\]
Lemma~6.1 bounds the left side below by
$K_0s^{-1}\tn{v'(k\alpha+l\beta)}
\gg_\alpha X^{1/3-1/20}/s$, a contradiction.
Taking $s=1$ gives the first conclusion, with $B>A+L$ for a test
of norm at most $(\log X)^L$.

For the twisted result assume a correlation larger than
$(\log X)^{-A}$. Apply quantitative Leibman to the polynomial
orbit $(g_0^{a+vt}\Gamma,\phi(t)\bmod1)$ on
$(H/\Gamma)\times\T$, with circle of filtration degree $d$
and product filtration degree $\max(2,d)$.
The test $(g\Gamma,z)\mapsto\Phi(g\Gamma)\ee(-z)$ has zero
Haar mean and polylogarithmic Lipschitz norm. The resulting
horizontal character $(k,l,m)$ has size $(\log X)^{O(1)}$ and
\[
 \norm{v(k\alpha+l\beta)t+m\phi(t)}_{C^\infty[K]}
 \ll(\log X)^{O(1)}.
\]
If $m=0$, Lemma~6.1 contradicts the inequality as in the first
part. For $m\ne0$ set $s=|m|$. Express the scalar polynomial
on the left as $b_0+\sum_{i=1}^d b_i\binom{t}{i}$.
Its smoothness bound gives integers $p_i$ with
$|b_i-p_i|\ll K^{-i}(\log X)^{O(1)}$.
Set $P(t)=\sum_i p_i\binom{t}{i}$ and
$E(t)=\sum_i(b_i-p_i)\binom{t}{i}$. For a constant $c$,
\[
 m\phi(t)=-(k\alpha+l\beta)(a+vt)+P(t)+E(t)+c .
\]
The factor $\ee(-P(t)/m)$ is constant on classes modulo $sd!$,
because $i!\binom{t}{i}$ is an integer polynomial and every
difference at step $sd!$ is divisible by $sd!$.
Also
$\sum_{t<K}|E(t+1)-E(t)|\ll(\log X)^{O(1)}$,
by the difference formula for binomial coefficients.

On the quotient define
\[
 \widetilde\Phi((x,y,z)\Gamma)
 =\Phi(D_s(x,y,z)\Gamma)\ee(\operatorname{sgn}(m)(kx+ly)).
\]
This absorbs the linear term on the dilated orbit. The map
$D_s$ has Lipschitz constant $O(s^2)$ in the fixed right-invariant
metric, so $\widetilde\Phi$ has polylogarithmic norm. The vertical
integral of $\widetilde\Phi$ is zero because $z\mapsto s^2z$ wraps the
vertical circle $s^2$ times and the horizontal phase is constant
along the fibre. For each residue class of $t$ modulo $sd!$,
the step of the dilated group orbit is determined by $v'=vsd!$,
with $v'\le X^{1/5000}$ and $r_j\nmid v'$.
These follow from the original $v$ range and $sd!<r_j$.
For a prefix of at least $X^{1/3}$, the first part bounds the
prefix sum of $\widetilde\Phi$ by its length times an arbitrary
inverse power of $\log X$. Shorter prefixes have trivial bound
$X^{1/3}(\log X)^{O(1)}$. A full class has
$\gg X^{1/2}(\log X)^{-O(1)}$ terms, so the short-prefix
bound also saves any inverse logarithm relative to that length.
This is the uniform prefix estimate (44), labelled
\code{e:phitildeonprog} in the manuscript.
Partial summation removes $\ee(-E(t)/m)$ at polylogarithmic cost.
Summing over the classes contradicts the assumed correlation.

Finally choose $\varepsilon=(\log X)^{-A-1}$. Let
$\chi_\varepsilon$ vanish at distance at most $\varepsilon$ from
the chart boundary $y=0\bmod1$, be $1$ from distance
$2\varepsilon$ on, and have Lipschitz norm $O(\varepsilon^{-1})$.
On the chart apply the twisted bound to
\[
 (x,y,z)\Gamma\longmapsto \chi_\varepsilon(\{y\})
     (F(z-x\lfloor y\rfloor)-\mu_F).
\]
The test extends Lipschitzly by zero near the boundary, has norm
$O_F(\varepsilon^{-1})$ and zero vertical mean by circle
translation. A Lipschitz majorant of the boundary strip, supported
within distance $3\varepsilon$ of $0$ with integral $O(\varepsilon)$,
shows by scalar equidistribution that removing the cutoff costs
$O_F(\varepsilon)$. This proves the centered bound uniformly.
\end{proof}
\noindent\textbf{Lean status.} Lemma~\ref{checked-group-model} checks
the filtered orbits, rational adapted Mal'cev bases, quotient metrics
and Haar measures. Lemma~\ref{checked-green-tao-adapter} applies the
general quantitative Leibman input on these two nilmanifolds and
proves the scalar and product criteria. These criteria give the weighted-prefix
bound by the factorial congruence, binomial rounding, the exact
Heisenberg covering identity, and variation control. Mean zero on the
cover follows by integration on the vertical fibres.
The uniform selection of scales from the dilated horizontal and scalar
estimates yields the twisted logarithmic saving for real polynomials
of degree at most $d$. The centered consequence (27) is checked
from the same two specialized character criteria. The Lean chart
test uses a piecewise linear cutoff. On each vertical fibre it has
mean zero by circle translation. Its right lattice invariance and
its Lipschitz estimate for the quotient distance are proved before
taking the orbit average. The scale $\varepsilon=(\log X)^{-A-1}$
balances the required test norm and the seam removal, whose lost
mass is controlled by the horizontal scalar estimate.

\begin{lemma}[4.5, multiplication by a localized vertical mode]\label{l:vertical-mode-multiplier}
\lean{NilBohr.localized_vertical_multiplier,NilBohr.localized_vertical_multiplier_of_bernstein}\leanok
\uses{selector,phase-tests}
Let $S\ge3$, let $I$ be a finite nonempty integer interval, and
let $a,v,h\in\Z$. If the Lipschitz function $\chi$ is supported on
an interval of length $<1/2$ inside $(0,1)$, then
\[
 \norm{f(t)\chi(\{\beta(a+vt)\})\ee(hZ(a+vt))}_{U^S[I]}
 \le\norm{\widehat\chi}_{\ell^1}\norm{f}_{U^S[I]} .
\]
\end{lemma}
\source{l:vertical-mode-multiplier}
\begin{proof}[Proof in the manuscript]
\proves{l:vertical-mode-multiplier}\uses{checked-cube,checked-norm-foundations}\leanok
Expand the conjugated products for a cube. If some
vertex lies outside $I$ or its cutoff vanishes, both products vanish.
Otherwise write $n_\omega=a+v(t+\omega\cdot u)$
and $b_\omega=\lfloor\beta n_\omega\rfloor$.
In each two-face,
$b_{00}+b_{11}-b_{01}-b_{10}$ is minus the corresponding
alternating sum of fractional parts. Its absolute value is
less than $2(1/2)=1$, because all part values lie in the
same short interval. Being an integer, it vanishes. Therefore
the increment in each direction is independent of the other
coordinates. Hence $b_\omega$ is affine on the cube.
Since $n_\omega$ is affine, $Z(n_\omega)$ is quadratic in the
bits of $\omega$. Every monomial omits at least one bit when $S\ge3$,
so pairing opposite vertices along that direction gives alternating
sum zero. Thus the two localized cube sums with and without the
vertical exponential agree exactly, giving equality of their
normalized $U^S$ norms.

For the remaining horizontal factor, Lipschitz regularity gives an
absolutely summable Fourier series. In fact Cauchy--Schwarz and
Parseval give
\[
 \sum_{k\ne0}|\widehat\chi(k)|
 \le\left(\sum_{k\ne0}k^{-2}\right)^{1/2}
    \left(\sum_{k\ne0}k^2|\widehat\chi(k)|^2\right)^{1/2}
 \ll\norm{\chi'}_2.
\]
The series converges uniformly, so expand
$\chi(\{\beta(a+vt)\})$ in the phases $\ee(k\beta(a+vt))$.
Each preserves the $U^S$ norm. The triangle inequality yields the
claimed multiplier bound.
\end{proof}
\noindent\textbf{Lean status.} The multiplier is checked under the
manuscript's hypotheses. Lemma~\ref{checked-norm-foundations} supplies
the triangle inequality and the absolute Fourier bound. The Fourier
limiting argument uses finite approximation and continuity of the cube
norm.

\section{The smooth-modulus prime model and count}

Throughout this section fix $c_0\in(0,1/11)$ and put
$\widetilde Q_N=\lceil\exp((\log N)^{c_0})\rceil$.
Write $\mathcal L(w)=\operatorname{lcm}(1,\ldots,w)$ and let
$P^+(q)$ be the greatest prime factor, with $P^+(1)=1$.
The smooth-modulus hypotheses in (17) are
\[
 w_0\le w\le(\log\widetilde Q_N)^{c_{\rm rel}},\qquad
 8\prod_{p\le w}p\mid q,\qquad P^+(q)\le w,\qquad
 \log q\le w^2.
\]

\begin{proposition}[4.1, high-order relative prime model]\label{p:high-order-relative-prime}
\lean{NilBohr.high_order_relative_prime_of_analytic_inputs,NilBohr.high_order_relative_prime_of_kmtt}\leanok
\uses{selector,phase-tests}
Fix $d\ge1$ and $S\ge2$. There exist
$c_{\rm rel},\eta_{\rm rel}>0$, $w_0\ge16$ such that, for all
large $N$, all smooth moduli above, and all integer $M$ with
$N^{1/d}\exp(-d\sqrt{\log N})\le M\le N^{1/d}$,
\[
 \norm{\Lambda_q-1}_{U^S[M]}\ll_{d,S}w^{-\eta_{\rm rel}}.
\]
In particular $\E_{y\in[M]}\Lambda_q(y)=1+O_{d,S}(w^{-\eta_{\rm rel}})$.
\end{proposition}
\source{p:high-order-relative-prime}
\begin{proof}[Proof in the manuscript]
\proves{p:high-order-relative-prime}\uses{external-inputs,checked-norm-foundations}\leanok
Banks's explicit bound gives $(1-\beta_0)\log r\ge0.003$
for a real zero of a primitive nonprincipal real character of
conductor $r\ge3$. For $r\le\widetilde Q_N$ this excludes a
KMTT exceptional zero, whose condition would require
$(1-\beta_0)\log\widetilde Q_N\le0.001$.
Thus the Siegel model in KMTT Lemma~3.1 is the
$\Cramer$ model at cutoff $\widetilde Q_N$.

Let $Y=qM+1$, for which $\log Y\asymp_d\log N$.
The cutoff range of KMTT Lemma~3.1 follows for all sufficiently
large $N$ as $c_0<1/11$. Apply it at ambient length $Y$.
Dilation to $y\in[M]$ with residue $1$ incurs at most
$q^{(S+1)/2^S}$, using KMTT Lemma~2.2 and the comparable
ambient intervals. Consequently the difference from the normalized
$\Cramer$ weight has norm
\[
 \ll_S q^{(S+1)/2^S}\exp(- (\log\widetilde Q_N)^{c_S}).
\]
Choose $2c_{\rm rel}<c_S$ to absorb the factor from $\log q\le w^2$.
The error is $O(w^{-1})$.
KMTT Lemma~4.3(i) with $b=1$ and full interval $[M]$
bounds the normalized $\Cramer$ weight minus $1$ by
$O(w^{-\eta_{\rm rel}})$ under exactly the smoothness and
scale constraints. Use the triangle inequality and normalized
Gowers monotonicity.
\end{proof}
\noindent\textbf{Lean status.} The independent premises are Banks's
zero exclusion and the two KMTT prime-model bounds. Lean supplies
dilation, the Gowers laws, cutoff compatibility and absorption of the
modulus loss.

\begin{lemma}[4.4, centered selector count]\label{l:unweighted-selector-relative-range}
\lean{NilBohr.unweightedSelectorComparison_of_local_phase}\leanok
\uses{selector,d:Cq}
Fix $C,d,k$, selector data and a Lipschitz $F$. For $N$ sufficiently
large and smooth moduli with $c_{\rm rel}$ sufficiently small, put
$M=\lfloor(N/q^{d-1})^{1/d}\rfloor$. For every fixed $A>0$,
the bound
$|T_{N,M,F_q-\mu_F}(\mathbf f)|\ll_{C,d,k,\alpha,\beta,F,A}w^{-A}$
holds uniformly for the distinct-degree $(C,q)$-families
and $1$-bounded inputs supported on $[N]$.
\end{lemma}
\source{l:unweighted-selector-relative-range}
\begin{proof}[Proof in the manuscript]
\proves{l:unweighted-selector-relative-range}\uses{t:main,p:heisenberg-phase}\leanok
Normalize the weight by $1+2\norm F_\infty$.
Apply Theorem~4.2 with $\rho=w^{-A-1}$.
Here $\log q\le w^2\ll(\log N)^{2c_0 c_{\rm rel}}$,
so the stated range for $q$ holds. The tests have steps $q^b$
and lengths $\gg Mq^{-O(1)}w^{-O(1)}=N^{1/d-o(1)}$.
Set $X=qM+1=N^{1/d+o(1)}$. Their $Z$-arguments form
progressions of length at least $X^{1/2}$ and step
$q^{b+1}\le X^{1/10000}$. Since $r_j>w$ and $P^+(q)\le w$,
the step is coprime to $r_j$. The affine change of test variable
preserves degree. Proposition~4.3 bounds each test by
$O_D((\log N)^{-D})$ for every fixed $D$.
Choose $D$ large enough after the exponent in Theorem~4.2 and
use $w\ll\log N$.
\end{proof}
\noindent\textbf{Lean status.} The deduction is checked from
Theorem~4.2 and the centered polynomial-phase estimate of
Proposition~4.3. The length and step restrictions, the normalization
of the test weight, and the affine pullback of the phase norm
are proved as part of the deduction.

\begin{theorem}[3.2, smooth-modulus count comparison]\label{t:relative-selector-comparison}
\lean{NilBohr.RelativeSelectorComparison,NilBohr.relativeSelectorComparison_of_obstructions,NilBohr.relativeSelectorComparison_of_centered_phase,NilBohr.relativeSelectorComparison_of_progression_estimates}\leanok
\uses{selector,d:Cq}
Fix $C,d,k$, selector data and Lipschitz $F:\T\to\C$.
For all sufficiently large $N$ and all smooth moduli above
with suitable constants $c_{\rm rel},\eta_{\rm sel}, w_0$,
take $M=\lfloor(N/q^{d-1})^{1/d}\rfloor$.
For $(C,q)$-shifts of distinct increasing degrees ending in $d$
and $1$-bounded $f_i$ supported on $[N]$,
\[
 T_{N,M,\Lambda_qF_q}(\mathbf f)
 =\mu_F T_{N,M,1}(\mathbf f)+O(w^{-\eta_{\rm sel}}).
\]
The constants are uniform in the shifts, $q,w,N$ and inputs.
\end{theorem}
\source{t:relative-selector-comparison}
\begin{proof}[Proof in the manuscript]
\proves{t:relative-selector-comparison}
\uses{l:unweighted-selector-relative-range,p:high-order-relative-prime,l:vertical-mode-multiplier,p:heisenberg-phase,external-inputs,checked-selector-cutoffs}\leanok
By linearity, subtracting the constant model gives weights
$(\Lambda_q-1)F_q$ and $F_q-\mu_F$. Lemma~4.4
gives $O(w^{-1})$ for the second. For the first, let a smooth
$0\le\chi\le1$ vanish outside $[\varepsilon,1-\varepsilon]$
in the horizontal coordinate and equal one on
$[2\varepsilon,1-2\varepsilon]$. Choose
$\norm\chi_{\rm Lip}\ll\varepsilon^{-1}$ and split it into
a fixed number of pieces $\chi_\nu$ of support length $<1/2$.
Expand $F_q$ by its absolutely convergent Fourier series.
KMTT Proposition~3.3 bounds each localized count by a
$U^S[M]$ norm at a fixed $S\ge3$. By Lemma~4.5 and the
Sobolev bound $\sum_\nu\norm{\widehat\chi_\nu}_{\ell^1}
\ll\varepsilon^{-1}$, the total localized bound is
$O_F(\varepsilon^{-1}\norm{\Lambda_q-1}_{U^S[M]})$.

For the boundary, use $|\Lambda_q-1|\le(\Lambda_q-1)+2$.
After absolute values, the desired bound reduces to the mean
of $1-\chi(\beta qy)$ and
$|\E_{y\in[M]}(\Lambda_q(y)-1)(1-\chi(\beta qy))|$.
Use Proposition~4.3 for the mean, with $a=0,v=q,X=qM+1$.
The first term is $O(\varepsilon+(\log X)^{-1})$ when
$\varepsilon^{-1}$ has at most polylogarithmic growth.
For the second, expand $1-\chi$ into its absolutely
summable linear phases and use Gowers monotonicity and
phase invariance for a bound $O(\varepsilon^{-1}
\norm{\Lambda_q-1}_{U^S[M]})$.

The scale of $M$ qualifies for Proposition~4.1 since
$\log q=o(\sqrt{\log N})$. Set
$\varepsilon=w^{-\eta_{\rm rel}/2}/32$, decreasing
$\eta_{\rm rel}$ to at most one. Then
$(\log X)^{-1}\ll\varepsilon$, and the combined bound
is $O(w^{-\eta_{\rm rel}/2})$.
\end{proof}
\noindent\textbf{Lean status.} The normalized finite inequalities for
the localized count, the boundary, and their sum are checked. The
quantitative cutoff construction uses three tent functions on the circle.
A positive floor on their common denominator gives a partition off the
boundary strip and a Lipschitz bound proportional to $\varepsilon^{-1}$.
The boundary integral follows from the length of a circle ball.
The uniform comparison uses Theorem~4.2, the high-order prime-model
estimate and KMTT's weighted counting estimate on a common sequence
of admissible denominators. The scalar and product specializations of
quantitative Leibman give the horizontal and centered tests. The
recurrence wrapper supplies its Gowers and Fourier estimates from
Lemma~\ref{checked-norm-foundations} and weak regularity from the
same inverse input used in Theorem~4.2.

\section{Density increment and recurrence}

\paragraph{Deriving the weak-regularity input.}
Write $k=\ell+1$, $u=q/\epsilon$ and
$\rho=\epsilon/(2k)$. With inverse witness bounds
$q'\le A\rho^{-B}$, $b'\le b_0$, and correlation at least
$a\rho^b N$, take the common step $q^J$ with $J=b_0+1$.
An increased resolution exponent makes every eligible $q'$ divide
the least common multiple in the hypothesis, hence divide $q$.
Thus $q'q^{b'}\mid q^J$. Put $E=B+J+b$ and take the coordinate
resolution $R_i=\lfloor c u^{-E}M^{\nu_i}\rfloor$ for small fixed
$c>0$, where $\nu_0=\deg P_1$ and $\nu_i=\deg P_i$ for $i>0$.
If any floor argument is less than two, the small-scale alternative
holds. Otherwise the bound on $M$ makes each nonempty residue
class large enough for this resolution. Divide each into atoms with
between $R_i$ and $2R_i$ elements. This yields both the common
step and the required lower bound on the atom lengths.

If replacing the functions by their atom means changes the count
by more than $\epsilon$, telescope one coordinate at a time and use
the residual divided by two in a term of size greater than $\rho$.
If the residual is at index zero, reindex the spatial sum by
$x\mapsto x+P_1(y)$ and swap the first two inputs. The new shifts are
$-P_1$ and $P_i-P_1$ for $i>1$. The coefficient class enlarges
from $C$ to $2C^2+C$, and all degrees stay in order. The supported
base functions justify both spatial cutoffs. The inverse theorem
now produces a test for the residual. On each atom subtract its
value at a fixed point, using the exact zero atom mean. The remaining
test is bounded by the inverse Lipschitz constant times the atom
diameter in units of the inverse step. Summed over the atoms, this
is at most $2A(2k)^B c\,u^{-b}N$. Choose $c$ so this is less than
$a\rho^b N$ since $u^{-b}\le\epsilon^b$. This contradicts the
inverse correlation. Lemma~\ref{checked-weak-regularity} verifies
this argument for the interface used in Proposition~3.1.

\begin{proposition}[3.1, one density increment]\label{p:selected-density-increment}
\lean{NilBohr.DensityIncrementFor,NilBohr.densityIncrementFor_of_relativeSelectorComparison}\leanok
\uses{selector,d:Cq}
There exist $c,\eta,K,E,C_1>0$ and $B_*\in\N$, depending only on the
fixed data of Theorem~1.2, with the following property for all large $N$.
Let $A\subset[N]$ have density $0<\delta\le1$ and suppose
\[
 K\delta^{-E}\le w\le(\log N)^c,\quad
 \mathcal L(w)\mid q,\quad P^+(q)\le w,\quad \log q\le w^2.
\]
Take $M=\lfloor(N/q^{d-1})^{1/d}\rfloor$.
If $A$ contains no pairwise distinct configuration with shifts
$P_{i,q}(y)$ and $qy+1\in\mathcal P_{Z,J}$, then there are
$1\le B\le B_*$, $a\in\Z$ and $N'\in\N$ with
\[
 a+q^B[N']\subset[N],\quad
 |A\cap(a+q^B[N'])|\ge(1+\eta)\delta N',\quad
 N'\gg q^{-C_1}\delta^{C_1}M^{\deg P_1}.
\]
\end{proposition}
\source{p:selected-density-increment}
\begin{proof}[Proof in the manuscript]
\proves{p:selected-density-increment}\leanok
\uses{t:relative-selector-comparison,checked-weak-regularity,checked-density}
Choose a nonnegative Lipschitz minorant $F\le\one_J$ with
$\mu_F>0$. For $w\ge16$ the divisibility
$\mathcal L(w)\mid q$ implies $8\prod_{p\le w}p\mid q$,
so Theorem~3.2 applies for sufficiently small $c$ and large $K$.
Write $k=\ell+1$. For a small fixed $\tau>0$ and
$1\le y\le\tau M$, every shift is at most $N/20$ in
absolute value. This follows from the coefficient bounds,
$q^{e-1}(\tau M)^e\le\tau q^{d-1}M^d\le\tau N$ for $e\ge1$.
There are at least $N/2$ admissible $x$ for each such $y$,
so
$T_1(\delta\one_{[N]},\ldots,\delta\one_{[N]})\ge c_3\delta^k$
with a fixed $c_3>0$.

For $\one_A$ inputs, the weighted count can receive mass only
from higher prime powers or a collision of shifts. Each polynomial
difference is nonzero even when one index is zero, so there are
$O_{d,k}(1)$ parameters producing a collision. The total
contribution is
\[
 O_{d,k}\left(\frac{\log(qM+2)}M+
      \frac{(qM+1)^{1/2}(\log(qM+2))^2}M\right)
 =O_{d,k}(N^{-1/(3d)}) .
\]
Here $M=N^{1/d-o(1)}$ and the higher-prime-power bound
follows by summing over exponents $j\ge2$ and using
$\Lambda(m)\le\log(qM+2)$. The error is $o(\delta^k)$
since $\delta^{-1}$ is polylogarithmic. Take $E\eta_{\rm sel}\ge k$
and then $K$ large, so the comparison error is at most
$\mu_F c_3\delta^k/4$. The unweighted count for $\one_A$
is then less than the constant-density count by at least
$c_3\delta^k/2$.

Apply KMTT Lemma~5.3 with accuracy
$\varepsilon=c_3\delta^k/8$. Increase $E$ and $K$ so the
required divisibility by $\mathcal L(\lceil\varepsilon^{-A'}\rceil)$
holds. The small-scale alternative $M\ll (q/\varepsilon)^{O(1)}$
is excluded because its right side is $N^{o(1)}$.
The factors have common modulus $q^B$ and atom lengths at
least a constant times $q^{-C_1}\delta^{C_1}M^{\deg P_1}$
because each resolution exponent is at least $\deg P_1$.
Let $H_i$ be their conditional expectations of $\one_A$,
extended by zero. Weak regularity changes the count by at most
$\varepsilon$, so the gap from the constant-density count is
at least $3c_3\delta^k/8$.

If all atoms have density $\le(1+\eta)\delta$, mass preservation gives
$\norm{H_i-\delta\one_{[N]}}_{\ell^1}\le2\eta\delta N$.
Replace one of the $k$ functions at a time. Translation in $x$
preserves the $\ell^1$ norm and the other factors contribute at most
$((1+\eta)\delta)^{k-1}$. The total change is at most
$2k\eta(1+\eta)^{k-1}\delta^k$, contradicting the gap for fixed
small $\eta$. Hence an atom supplies the increment.
\end{proof}
\noindent\textbf{Lean status.} The uniform density increment is
checked from the count comparison of Theorem~3.2 and KMTT weak
regularity. Lean fixes the constants before ranging over all admissible
scales, moduli, densities and input sets. The singular parameters
with coincident shifts are counted separately.

\begin{theorem}[1.2, quantitative selected-prime recurrence]\label{t:recurrence}
\lean{NilBohr.HeisenbergRecurrence,NilBohr.heisenbergRecurrence_of_analytic_inputs,NilBohr.heisenbergRecurrence_of_obstructions,NilBohr.heisenbergRecurrence_of_centered_phase,NilBohr.heisenbergRecurrence_of_relativeSelectorComparison}\leanok
\uses{selector,d:Cq}
For each fixed polynomial family and selector data there is
$\kappa>0$ such that for every sufficiently large $N$, every
$A\subset[N]$ with $|A|\ge N(\log N)^{-\kappa}$ contains a
pairwise distinct configuration with prime parameter
$p\in\mathcal P_{Z,J}$, so $n=p-1$.
\end{theorem}
\source{t:recurrence}
\begin{proof}[Proof in the manuscript]
\proves{t:recurrence}\uses{p:selected-density-increment,checked-density}\leanok
Assume the original set contains no selected configuration.
Fix $w=\lceil K\delta^{-E_1}\rceil$ with $E_1\ge E$ chosen below,
$q_0=\mathcal L(w)$, and rescale a densest nonempty residue class.
The new set has density $\delta_0\ge\delta$ and length
$N_0\ge N/(2q_0)$ for large $N$ since $\log q_0\ll w$.
The affine lift proves it has no selected configuration with the
rescaled polynomials and the same prime selector.

At stage $i$, each application of Proposition~3.1 multiplies the
density by at least $1+\eta$ and produces $q_{i+1}=q_i^{B_i+1}$.
All moduli remain powers of $q_0$, so divisibility and smoothness
persist, and increasing density preserves the lower bound on $w$.
The affine identities in the rescaling definition preserve the
absence of a selected configuration. Let
$i_*=\lfloor\log(\delta^{-1})/\log(1+\eta)\rfloor+1$.
For $i\le i_*$,
\[
 \log q_i\le (B_*+1)^i\log q_0
       \ll\delta^{-D}w
\]
with fixed $D$. Choose $E_1>D$ and increase $K$ to bound
this by $w^2$. Put $r=\deg P_1$. The lower length in the
increment and the floor defining $M$ give
\[
 \log N_{i+1}\ge(r/d)\log N_i
      -O(\log q_i+\log(\delta^{-1})+1)
\]
when the current stage is in range. With $i\le i_*$,
$(r/d)^i\gg\delta^{D'}$ for fixed $D'\ge0$. Iteration gives
$\log N_i\ge c_2\delta^{D'}\log N
-O((1+\log(\delta^{-1}))(w^2+\log(\delta^{-1})))$.
For $\delta\ge(\log N)^{-\kappa}$ choose sufficiently small
$\kappa$ after all previous constants. This yields
$\log N_i\ge(\log N)^{3/4}$ and
$w\le(\log N_i)^c$ at every stage up to $i_*$.
The size conditions therefore persist by induction.
After $i_*$ steps the density exceeds one, a contradiction.
\end{proof}
\noindent\textbf{Lean status.} The iteration and the density increment
are checked, including the asymptotic length and modulus estimates.
The composed deduction uses the literature interfaces for the prime
models, zero exclusion, local inverse, weighted count and quantitative
Leibman. Lemma~\ref{checked-weak-regularity} supplies weak regularity.

\begin{corollary}[1.3, bracket recurrence]\label{c:bracket-recurrence}
\uses{selector}\lean{NilBohr.bracket_recurrence_of_heisenberg_recurrence,NilBohr.bracketRecurrence_of_analytic_inputs,NilBohr.bracketRecurrence_of_obstructions,NilBohr.bracketRecurrence_of_centered_phase}\leanok
Theorem~1.2 holds with $Z_{\rm br}$ in place of $Z$.
\end{corollary}
\source{c:bracket-recurrence}
\begin{proof}[Proof in the manuscript]
\proves{c:bracket-recurrence}\uses{t:recurrence,checked-floor}\leanok
Choose $J_0$ with $-\overline J_0\subset J$ and choose
$\varepsilon<\min(\beta,1-\beta)$ with
$(3+|\alpha|)\varepsilon<\operatorname{dist}
(-\overline J_0,\T\setminus J)$.
Enlarge the family by $n$ and $n^2$ only when their degrees are
absent. Write its linear and quadratic shifts as $an$
and $bn^2+cn$ with $ab\ne0$. Partition $A$ into finitely many
torus boxes of side length less than $\varepsilon$ for the coordinates
\[
 x\longmapsto
 \left(\frac{\alpha x}a,\frac{\beta x}a,
       \frac{\theta(1-c/b)x}a,\frac{\theta x}b\right)\bmod1.
\]
Apply Theorem~1.2 to a largest cell with the enlarged family
and $J_0$. The fixed partition loss is absorbed by decreasing
$\kappa$. The differences for the linear shift bound
$\tn{\alpha n}$, $\tn{\beta n}$, and
$\tn{\theta(1-c/b)n}$ by $\varepsilon$. The quadratic shift
bounds $\tn{\theta n^2+(\theta c/b)n}$ by $\varepsilon$.
Thus $\tn{\theta(n^2+n)}<2\varepsilon$.

Let $m$ be a nearest integer to $\beta n$. The choice of
$\varepsilon$ forces $\lfloor\beta(n+1)\rfloor=m$ and
$m-\lfloor\beta n\rfloor\in\{0,1\}$. Direct substitution gives
\[
 Z_{\rm br}(n)+Z(n)=\theta(n^2+n)
     +\alpha n(m-\lfloor\beta n\rfloor)-\alpha(\beta n-m).
\]
Its torus norm is less than $(3+|\alpha|)\varepsilon$.
Since $Z(n)\bmod1\in J_0$, the bracket phase belongs to $J$.
Retaining the original shifts preserves pairwise distinctness.
\end{proof}
\noindent\textbf{Lean status.} Checked from Theorem~1.2 as an
explicit premise.
Lean fills all missing positive degrees up to $\max(d,2)$ to
index the enlarged family by degree, and reflects the central
half of $J$. Restricting the configuration to the original
shifts gives the stated conclusion.

The abstract takes $P_i(n)=n^i$, $\alpha=\sqrt3$ and
$J=(1/4,1/2)$. If $u$ is the nearest integer to $m\sqrt3$
with $m\ge1$, then $|3m^2-u^2|\ge1$ and
$m\sqrt3+u<4m$. Consequently $m\tn{m\sqrt3}\ge1/4$,
which verifies the stated bad-approximability hypothesis in that
example. The bracket recurrence yields precisely the selector
$\{\sqrt3\,p\lfloor\beta p\rfloor\}\in J$.

\begin{corollary}[1.1, finite Bohr constraints]\label{c:bohr}
\uses{d:Cq}\lean{NilBohr.bohr_recurrence}\leanok
If the polynomial family includes $P_1(n)=an$ with $a\ne0$,
one can force finitely many conditions
$\tn{\alpha_j(p-1)}<\varepsilon$ in the shifted-prime recurrence
conclusion, with a fixed polylogarithmic density threshold.
\end{corollary}
\source{c:bohr}
\begin{proof}[Proof in the manuscript]
\proves{c:bohr}\uses{external-inputs}\leanok
Partition $A$ by a fixed torus grid of side length
$<\varepsilon$ for the coordinates $(\alpha_jx/a)_j$.
A largest cell has a fixed proportion of $|A|$. Apply KMTT
Theorem~1.2 to that cell, decreasing the density exponent to
absorb the fixed loss. Both $x$ and $x+a(p-1)$ lie in the same
cell, so their coordinate differences give the conditions.
\end{proof}
\noindent\textbf{Lean status.} Checked from the KMTT
shifted-prime bound. The signed coefficient, strict inequalities, and all original shifts
are retained.

\section{Prime-weighted phase tests}

\begin{proposition}[6.2, prime-weighted local cancellation]\label{p:prime-phase}
\lean{NilBohr.PrimeLocalPhase,NilBohr.primeLocalPhase_of_literature,NilBohr.heisenbergPrimeLocalPhase_of_literature}\leanok
\uses{selector,phase-tests,d:Cq}
Fix $d\ge1$, $A,B,K,L>0$ and Lipschitz $F$. For large $N$,
let $1\le q\le(\log N)^L$, $M=\lfloor(N/q^{d-1})^{1/d}\rfloor$,
and $P=I\cap(u+q^b\Z)$ with $I\subset[M]$, $|I|\ge M(\log N)^{-K}$,
$0\le b\le B$ and $b,u\in\Z$. Then
\[
 \sup_{\deg\phi\le d}\abs{\E_{y\in P}\Lambda_q(y)
 (F_q(y)-\mu_F)\ee(-\phi(y))}
 \ll_{d,\alpha,\beta,F,A,B,K,L}(\log N)^{-A}.
\]
\end{proposition}
\source{p:prime-phase}
\begin{proof}[Proof architecture]
\proves{p:prime-phase}
\uses{p:heisenberg-phase,l:heisenberg-horizontal,l:vertical-mode-multiplier,external-inputs,checked-prime-cramer,checked-prime-seam,checked-fourier-restriction}\leanok
Set $X=qM+1$. The substitution $m=qy+1$ maps $P$ to consecutive
points $P'$ of the reduced class $qu+1\bmod q^{b+1}$. For large $N$,
\[
 |P'|=|P|\ge |I|/q^b-1\gg X(\log X)^{-K-(B+1)L},
 \qquad
 \frac{\log N}{2d}\le\log X\le2\log N.
\]
In particular $P'$ is nonempty. The test becomes
$p(m)=\phi((m-1)/q)$ and keeps degree at most $d$.

\emph{Remove the horizontal seam.}
Take $\eta=(\log X)^{-A-2}$ and
$\chi_\eta(y)=\min(1,\max(0,(\tn y-\eta)/\eta))$.
It vanishes at distance at most $\eta$ from $0$ and equals
one from distance $2\eta$ on. The complement has Haar mean
$O(\eta)$ and Lipschitz norm at most $\eta^{-1}$.
Apply scalar equidistribution to the orbit with $m-1$
along $P'$. Its step $q^{b+1}$ is polylogarithmic and less
than $r_j$, and its length is at least $X^{1/2}$. The bound
$\Lambda(m)\le\log X$ shows that the cutoff costs
$O_F(\eta\log X)$, uniformly over all polynomial phases.

On the quotient, use the test
\[
 \Phi_\eta(g)=
  \chi_\eta(y(g)\bmod1)\bigl(F(z_{\rm red}(g)\bmod1)-\mu_F\bigr).
\]
Its value at $g(m-1)$ equals
$\chi_\eta(\beta(m-1)\bmod1)(F(Z(m-1))-\mu_F)$.
It has vertical mean zero, bounded sup norm, and quotient Lipschitz
norm $O_F(1+\eta^{-1})$, by the centered chart estimate in
Proposition~4.3. Set $f_\eta(m)=\Phi_\eta(g(m-1))$.

\emph{Replace the prime weight by the model.}
Let $S=\max(3,d+1)$ and
$Q=\lceil\exp((\log X)^{c_0})\rceil$.
Banks's real-zero exclusion and KMTT Lemma~3.1 at scale $X$ give
\[
 \norm{\Lambda-\Lambda_{\Cramer,Q}}_{U^S[X]}
 \ll_S \exp(-c(\log X)^{c_0c_S})
\]
for some $c>0$. First test a single vertical frequency $h\in\Z$ with a
horizontal multiplier $\chi$ supported in a chart interval shorter
than $1/2$. Expand $\one_{P'}$ using the congruence indicator, of
Fourier $\ell^1$ mass one, and the spanning interval indicator, of
Fourier $\ell^1$ mass at most $2+\log X$ on the group of size
$2X+1$. Every frequency introduced by these expansions is linear.
The mean bound for $U^S[X]$ and invariance under phases of degree
$<S$ control the test, including the degree-$d$ twist. Lemma~4.5
then bounds the localized mode by
\[
 C_S\frac{X}{|P'|}(2+\log X)
       \norm{\widehat\chi}_{\ell^1}
       \norm{\Lambda-\Lambda_{\Cramer,Q}}_{U^S[X]} .
\]
The bound is uniform in $h$. Bernstein's theorem gives absolute
Fourier convergence of $F-\mu_F$, so linearity of the finite
correlation and the triangle inequality extend the bound to its
full Fourier series with a factor $\norm{\widehat{F-\mu_F}}_{\ell^1}$.
An exact decomposition of $\chi_\eta$ into three short-supported
pieces follows by multiplying the horizontal chart partition by
$\chi_\eta$. The partition's boundary piece vanishes in this
product. Each piece has sup and Lipschitz norms bounded by
$O(\eta^{-2})$, so its Fourier mass satisfies the same bound.
The losses are all polylogarithmic, including $X/|P'|$, and the
ambient Gowers estimate absorbs every fixed logarithmic power.
This use of absolute convergence avoids truncating the vertical
Fourier series.

\emph{Sieve the model and keep the full centered chart test.}
Put $W=\prod_{p\le Q,\ p\nmid q}p$, $H=|P'|$, and
$\mathcal D=\lfloor X^{1/20000}\rfloor$. The model on $P'$ is
the coprimality indicator $\one_{(m,W)=1}$ times the Euler product
$\prod_{p\le Q}(1-1/p)^{-1}=O(\log X)$ by Mertens. The
dimension-one interval bound for these primes follows by removing
factors from the full product. For the sieve input use
Koukoulopoulos, \emph{The Distribution of Prime Numbers},
Theorems~3.4(c) and 19.1, in dimension one with local
densities $1/p$. Its prime set $p\le Q$ is the one obtained
by cutting at $Q+1$ with the manuscript's strict-cutoff convention.
The level ratio $U=\log\mathcal D/\log Q$ has
$U\gg(\log X)^{1-c_0}$. For $U\ge e^2$, the relative error
$U^{-U/2}\le e^{-U}$, and the Euler product in that main term has
size at most one. Thus the fundamental lemma supplies real
coefficients $|\lambda_\ell^\pm|\le1$ on
$\ell\mid W$, $\ell\le\mathcal D$ whose sums bracket
$\one_{(m,W)=1}$ and whose reciprocal-weighted gap is at most
$(\log X)^{-D}$ for every fixed $D>0$ and large $X$.
If no primes remain in $W$, the two envelopes supported at
divisor $1$ coincide with the indicator and have gap zero.

Write $H_\ell=|\{m\in P':\ell\mid m\}|$. The divisors are
coprime to $q^{b+1}$, so each intersection has step
$q^{b+1}\ell$ and $H_\ell=H/\ell+O(1)$. The midpoint of
the two envelopes incurs average absolute error
$O((\log X)^{-D}+\mathcal D/H)$ because its pointwise
error is at most half the envelope gap. Taking $D=A'+2$ gives
the saving in the manuscript's midpoint estimate. This controls the
weighted error after multiplying by the fixed bound for $|f_\eta|$.

For $r_j\nmid\ell$, the step $q^{b+1}\ell$ is coprime to
$r_j$ and at most $X^{1/10000}$. The length is at least $X^{1/2}$.
Apply the twisted orbit estimate to $\Phi_\eta$ along the
divisor progression, with the polynomial $p$ evaluated in its
affine coordinate. Its degree is still at most $d$, and
the full vertically centered test has polylogarithmic norm.
Relative to the common denominator $H$, each contribution is
$(H_\ell/H)O_D((\log X)^{-D})$. Sum the good divisors at
harmonic cost $O(\log X)$ and use the trivial bound on the
bad divisors, whose total contribution is
$O(\log X/r_j+\mathcal D/H)$. The ratio $\mathcal D/H$
saves every fixed logarithmic power, as does $r_j^{-1}$.
Choose $D$ to cover the harmonic loss and the Euler
normalization. The sieve thus bounds the Cram\'er contribution
without dividing $F-\mu_F$ into individual vertical modes.

The seam, model replacement, and Cram\'er terms add before
taking the supremum. Changing back to $y$ costs only
$\varphi(q)/q\le1$, and $\log X\ge (\log N)/(2d)$
converts the resulting bound to the stated $(\log N)^{-A}$
uniformly in $I,b,u,\phi$.
\end{proof}
\noindent\textbf{Lean status.} The checked bundle constructor uses the
ambient prime-model and zero-free inputs, quantitative Leibman, the Mertens
interval estimate, and the dimension-one fundamental lemma.
It supplies the normalized Gowers, Fourier and obstruction fields.
The reciprocal-square tail bound gives a uniform Fourier truncation
error at $|h|\le\eta^{-2}$ of at most $2\operatorname{Lip}(F)\eta$.
With $\Lambda(m)\le\log X$, this verifies the truncation error of the
printed argument.
Lemmas~\ref{checked-prime-seam},
\ref{checked-fourier-restriction}, and \ref{checked-prime-cramer}
give the three error bounds before the phase supremum. The main
theorem chooses its constants and threshold before quantifying over
$q,b,I,u$ and the test phase. The integer bound $B$ in Lean covers
real $B$ by rounding up. The uniform wrapper extracts the sequence
of prime approximants from the manuscript's admissibility hypothesis.


\section{Pointwise ergodic comparison}

\paragraph{Prime distribution from the available inputs.}
For the cutoff $Q=\lceil\exp((\log N)^{c_0})\rceil$, take sieve
level $D=\lfloor N^{1/20000}\rfloor$ and absolute gap
$\epsilon=(\log N)^{-2}$. If $W$ is the product of the sieving
primes, average both pointwise bounds on a complete period. A divisor
$d\mid W$ occurs with frequency $1/d$ there, and the sifted indicator
has mean $\varphi(W)/W$. On an actual positive interval of length $N$,
the corresponding counts differ from $N/d$ by at most one for each
divisor. The same bound holds on an affine progression with step
coprime to $W$, replacing $N$ by the number of its terms. The full
Cram\'er normalizer has size $O(\log N)$ by Mertens. Thus on $[N]$ the
model mean tends to one, with error $O(\log N(\epsilon+D/N))$.
For a fixed reduced class modulo $q$, omit its prime factors in the
sieve, then use $Q\ge q$ for large $N$. The full Euler normalizer times
the retained sieve density is exactly $q/\varphi(q)$. The class has
$N/q+O(1)$ terms, so the ambiently normalized model mean tends to
$1/\varphi(q)$. The checked residue Fourier projection contracts the
order-two norm of the model discrepancy. Its mean bound then passes
both model limits to Mangoldt. Removing prime powers and logarithmic
weights gives the two prime count asymptotics.

\paragraph{The irrational linear test.}
Fix irrational $\alpha$ and put $R=\lfloor\sqrt N\rfloor$. For each
large $N$ take a reduced $a/q$ with $1\le q\le R$ and
$\vartheta=\alpha-a/q$ satisfying $|\vartheta|\le1/(qR)$.
For any fixed denominator bound the finitely many small denominators
have a positive minimum of $\|q\alpha\|_{\mathbb T}$. As
$R\to\infty$, the bound $\|q\alpha\|_{\mathbb T}\le1/R$ therefore
forces $q$ past every fixed bound. This observation permits arbitrary
irrational $\alpha$ at the full sequence of cutoffs.

Replace the indicator in the Cram\'er correlation by the real
midpoint of its two sieve bounds. Against a complex test of norm
at most one, the error is bounded by the normalized half gap plus
$O((W/\varphi(W))D/N)$, since the expected divisor counts have
errors at most one. For $d\le D$ with $q\nmid d$,
\[
 \|d\alpha\|_{\mathbb T}\ge 1/q-d/(qR)\ge 1/(2q)
\]
eventually, by $2D\le R$. The geometric sum of $\ee(\alpha n)$
over the $d$-multiples in $[N]$ has norm at most $2q$. If $q\mid d$,
its phase equals $\ee(\vartheta n)$ on those multiples. Compare
the sum with $d^{-1}\sum_{n\in[N]}\ee(\vartheta n)$ by using the
endpoint of each complete block of length $d$. A fixed Lipschitz
bound for $\ee$ gives error $O(|\vartheta|N)$, and the last partial
block contributes at most one. After summing the divisor errors
and using $qR\le N$, the normalized error is
$O((W/\varphi(W))(D/N+D/R))$.

The remaining multiplier is the midpoint's rational Fourier
coefficient over a common period of $W$ and $q$. The same
bracketing compares it with the true normalized coprime Fourier
coefficient, now with no endpoint error. M\"obius inversion expands
the coprimality indicator. A geometric sum over the $d$-multiples
of the complete period vanishes unless $q\mid d$. If $q\nmid W$
the coefficient is zero. Otherwise it is $\mu(q)/\varphi(q)$ by
multiplicativity and the totient identity. Since $W$ is squarefree,
in this case $q\le2\varphi(q)^2$ by $p\le(p-1)^2$ for $p\ge3$.
The bound $\sqrt{2/q}$ follows. The remaining errors tend to zero
because $(W/\varphi(W))\ll\log N$, $D\le N^{1/20000}$ and
$R\asymp\sqrt N$. The ambient $U^2$ discrepancy is small against
every linear phase by the checked phase-norm comparison. Convert
the Mangoldt correlation to ordinary primes and use trigonometric
density for the continuous tests. Lemma~\ref{checked-prime-distribution}
records the three checked distribution interfaces.

Let $(X,\nu,T)$ be an invertible probability
measure-preserving system. For bounded measurable representatives
$f_i\in L^\infty(X)$, write $z_n(x)=\prod_{i=1}^\ell f_i(T^{P_i(n)}x)$.
On a common full-measure set the pointwise bound
$|z_n(x)|\le\prod_i\norm{f_i}_\infty$ holds for all relevant $n$,
by measure preservation and the countable family of iterates.
Define
\[
 \mathcal A_H^F(x)=\frac1H\sum_{n\in[H]}
       \Lambda(n+1)F(Z(n))z_n(x),\qquad
 A_H(x)=\frac1{\pi(H)}\sum_{p\le H}z_{p-1}(x).
\]

\begin{proposition}[7.1, smooth ergodic comparison]\label{p:ergodic-comparison}
\lean{NilBohr.smoothErgodicComparison_of_prime_model_distinct_degrees,NilBohr.smoothErgodicComparison_of_literature,NilBohr.smoothErgodicComparison_of_primeLocalPhase,NilBohr.smoothErgodicLogDecay_of_primeLocalPhase}\leanok
\uses{selector}
Fix the polynomial family with distinct positive degrees in arbitrary
order and zero constant terms. Fix the selector data, a Lipschitz
$F:\T\to\C$, and $A>0$. For every integer $H\ge2$,
\[
 \norm{\mathcal A_H^F-\mu_F\mathcal A_H^1}_{L^1(X)}
 \ll_{\mathbf P,\alpha,\beta,F,A}
    (\log H)^{-A}\prod_i\norm{f_i}_\infty.
\]
The difference tends to zero almost everywhere. Furthermore,
\[
 \frac1{\pi(H)}\sum_{p\le H}
   (F(Z(p-1))-\mu_F)z_{p-1}(x)\longrightarrow0
 \quad\text{for almost every }x .
\]
\end{proposition}
\source{p:ergodic-comparison}
\begin{proof}[Proof in the manuscript]
\proves{p:ergodic-comparison}\uses{p:prime-phase,external-inputs,checked-ergodic,checked-orbit,checked-prime-distribution}\leanok
Order the indices by degree and permute the functions with the
polynomials. Both the orbit product and the product of essential
supremum norms are invariant under this permutation. This chooses
the order before the probability system and the inputs are given.
By homogeneity use $1$-bounded inputs, treating a zero norm
by the common full-measure set.
Set $\vartheta(n)=\Lambda(n+1)(F(Z(n))-\mu_F)$ and
$L_H=2(1+\norm F_\infty)\log(H+2)$.
Proposition~6.2 gives
$\norm{\vartheta}_{u^{d+1}[H]}\ll_D(\log H)^{-D}$ for any
fixed $D>0$, with $q=1$ and $N=H^d$.
The weight $\vartheta/L_H$ is $1$-bounded.

Pair the centered difference with a $1$-bounded test $g$.
Use $x=T^mz$ and average this measure-preserving change of
variables over $m\in[H^d]$. The orbit functions are
$g_{0,z}(m)=g(T^mz)\one_{[H^d]}(m)$ and
$g_{i,z}(m)=f_i(T^mz)\one_{[-CH^d,CH^d]}(m)$
for a fixed $C$ containing all shifted supports. This writes
the pairing as the integral over $z$ of $L_H$ times the
discrete count with denominator $H^{d+1}$ and weight
$\vartheta/L_H$. Ter\"av\"ainen's Theorem~4.1 applies at
these exact scales. By $L^1$ duality the $L^1$ difference is
\[
 \ll_{\mathbf P} L_H
    (H^{-1}+L_H^{-1}\norm{\vartheta}_{u^{d+1}[H]})^{1/K_0}
\]
for some $K_0\ge1$. Take $D>K_0(A+1)$; bounded $H$ can be
absorbed by enlarging the constant.

For $H_j=\lfloor\exp(j^{1/3})\rfloor$ and $A>3$ the
$L^1$ errors are summable in $j$. Thus the sequence of
differences tends to zero almost everywhere along $H_j$.
For $H_j\le H\le H_{j+1}$ the same summands occur at both
scales and their magnitudes are $O_F(\log(H_{j+1}+2))$.
The normalized difference between the two scales is
\[
 O_F\left(\frac{H_{j+1}-H_j}{H_j}\log(H_{j+1}+2)\right)
 =O_F(j^{-1/3}).
\]
Here the increment of $j^{1/3}$ is $O(j^{-2/3})$ and the floor
contributes a negligible rounding error. The gap bound extends
convergence to all $H$ on the same full-measure set.

Write $a_p(x)=(F(Z(p-1))-\mu_F)z_{p-1}(x)$.
Drop the endpoint at $H+1$ at cost $O_F(\log H)$, and the
higher prime powers at cost $O_F(H^{1/2}(\log H)^2)$.
It follows that $B(H)=\sum_{p\le H}(\log p)a_p(x)=o(H)$.
Partial summation gives
\[
 \sum_{p\le H}a_p(x)
 =\frac{B(H)}{\log H}
    +\int_2^H\frac{B(t)}{t(\log t)^2}\,dt=o(H/\log H).
\]
The integral has that asymptotic, for instance by cutting at
$\sqrt H$ and using $B(t)=o(t)$ on the upper range.
The prime number theorem converts the estimate to the stated
normalized convergence.
\end{proof}
\noindent\textbf{Lean status.} The $L^1$ estimate holds in Lean for
every $A>0$ from the literature inputs of Proposition~6.2 and the
fixed-family Ter\'av\'ainen estimate. This includes the finite orbit
window, scaling for bounded inputs, and extraction of the $q=1$
test. The passage to almost-everywhere prime averages is checked
from exponent $6$ and the stated PNT asymptotics.

\begin{theorem}[1.4, Heisenberg selected prime averages]\label{t:pointwise}
\lean{NilBohr.PrimePointwiseFor,NilBohr.heisenberg_pointwise_of_prime_model,NilBohr.heisenberg_pointwise_of_analytic_inputs,NilBohr.heisenberg_pointwise_of_literature,NilBohr.primePointwiseFor_of_primeLocalPhase}\leanok
\uses{selector}
Under the fixed hypotheses and for every bounded input family,
$\pi_{Z,J}(H)/\pi(H)\to|J|$ and
\[
 \frac1{\pi_{Z,J}(H)}
   \sum_{p\le H,\ p\in\mathcal P_{Z,J}}z_{p-1}(x)
 \longrightarrow \lim_{H\to\infty} A_H(x)
 \quad\text{for almost every }x ,
\]
where the unrestricted limit exists almost everywhere.
\end{theorem}
\source{t:pointwise}
\begin{proof}[Proof in the manuscript]
\proves{t:pointwise}\uses{p:ergodic-comparison,external-inputs,checked-ergodic,checked-prime-distribution}\leanok
Apply Mirek--Wan--Wright Theorem~1.1(ii) to the polynomials
$P_i(t-1)$, whose degrees remain distinct, and inputs in
$L^{2\ell}(X)$. This gives pointwise convergence of
$H^{-1}\sum_{n\le H}\Lambda(n)z_{n-1}(x)$.
Drop higher prime powers and use partial summation and the prime
number theorem, now allowing the limiting function as the main
term, to obtain a limit $A_H(x)\to L(x)$.

For integer $j\ge1$ choose real Lipschitz functions
$0\le F_j^-\le\one_J\le F_j^+\le1$ and
$\mu_{F_j^+-F_j^-}<2^{-j}$. Such functions can be made
by linear ramps near the two boundary points, with
arbitrarily small ramp widths. Apply Proposition~7.1 to $F_j^-$
with the given inputs and to $F_j^+-F_j^-$ with unit inputs.
Intersect the resulting countable family of full-measure sets.
With normalization $1/\pi(H)$, the norm of the replacement
error from $F_j^-$ to $\one_J$ is bounded by
\[
 \prod_i\norm{f_i}_\infty
  \frac1{\pi(H)}\sum_{p\le H}(F_j^+-F_j^-)(Z(p-1)).
\]
The scalar limit bounds its limsup by
$2^{-j}\prod_i\norm{f_i}_\infty$. Also
$|\mu_{F_j^-}-|J||<2^{-j}$. Sending $j\to\infty$ gives
the comparison of the selected sum normalized by $\pi(H)$
with $|J|A_H(x)$, with error tending to zero. With all
inputs one, this gives the scalar count ratio.
Division by the positive limiting ratio proves the result.
All normalized averages are taken once the denominators are
positive, which the scalar limit implies eventually.
\end{proof}
\noindent\textbf{Lean status.} The conclusion is checked from
the literature inputs of Proposition~6.2, Ter\'av\'ainen's fixed-family
estimate, and Mirek--Wan--Wright. Lemma~\ref{checked-prime-distribution}
supplies the prime count asymptotic from the same prime-model and
sieve premises. The quantitative rate of Proposition~7.1 is an
intermediate checked deduction.

\begin{corollary}[1.5, bracket selected prime averages]\label{c:bracket-pointwise}
\lean{NilBohr.bracket_pointwise_of_prime_model,NilBohr.bracket_pointwise_of_analytic_inputs,NilBohr.bracket_pointwise_of_literature,NilBohr.bracket_pointwise_of_primeLocalPhase}\leanok
\uses{selector}
The conclusion of Theorem~1.4 holds for $Z_{\rm br}$ in place of $Z$.
\end{corollary}
\source{c:bracket-pointwise}
\begin{proof}[Proof in the manuscript]
\proves{c:bracket-pointwise}\uses{t:pointwise,checked-prime-distribution,checked-rotations}\leanok
First fix $h\in\Z\setminus\{0\}$ and a real polynomial
$\psi$ of degree at most two. Adjoin shifts of degrees one
and two if missing and call the increasing family $Q_i$.
The corresponding $g_i$ is the original input at $Q_i$ or $1$
for an added shift. Write
$\psi(n)=\psi(0)+\omega_1Q_1(n)+\omega_2Q_2(n)$
using the nonzero leading coefficients. For
$S(x,u,v)=(Tx,u+\omega_1,v+\omega_2)$ on $X\times\T^2$,
put $G_1=g_1(x)\ee(u)$, $G_2=g_2(x)\ee(v)$ and
$G_i=g_i(x)$ for $i\ge3$. The product of iterated
$G_i$ is $\ee(u+v-\psi(0))\ee(\psi(n))z_n(x)$.
Theorem~1.4 on this system and Fubini give for each fixed
open arc $K$,
\[
 \frac1{\pi(H)}\sum_{p\le H}
 (\one_K(Z(n))-|K|)\ee(\psi(n))z_n(x)\longrightarrow0,
 \qquad n=p-1,
\]
almost everywhere. Endpoints have zero relative density:
a centred arc around any endpoint has prime proportion tending
to the arc length by the scalar theorem, which can be arbitrarily
small. Thus the formula holds for half-open arcs as well.
Approximate $\ee(-ht)$ by its values on equal half-open
intervals of width $1/m$, for integers $m>|h|$.
Their constant contribution cancels since the roots of unity
sum to zero. The uniform error is at most $2\pi|h|/m$.
On a common full-measure set, send $m\to\infty$ to conclude
that $\ee(-hZ(n)+\psi(n))z_n(x)$ has prime average tending to zero.

Next put $t=\{\beta n\}$ and $c(t)=\one_{[1-\beta,1)}(t)$.
The exact carry identity is
\[
 Z_{\rm br}(n)+Z(n)=\theta(n^2+n)-\alpha t+\alpha(n+1)c(t).
\]
Exponentiating $h$ times gives
\begin{align*}
 \ee(h Z_{\rm br}(n))
 &=\ee(-hZ(n)+h\theta(n^2+n))\bigl[U_0(t)+\ee(h\alpha n)U_1(t)\bigr],\\
 U_0(t)&=\one_{[0,1-\beta)}(t)\ee(-h\alpha t),\qquad
 U_1(t)=\one_{[1-\beta,1)}(t)\ee(h\alpha(1-t)).
\end{align*}
The only possible discontinuities are at $0$ and $1-\beta$.
Approximate $U_0,U_1$ by continuous functions of modulus at
most one that agree outside $\delta$-neighborhoods of those
points. This can be done by linear interpolation in the complex
unit disk. Since $\beta$ is irrational, prime equidistribution
of $\{\beta(p-1)\}$ bounds the limiting proportion of discarded
terms by $4\delta$, hence bounds the limsup of the weighted
error by $16\delta\prod_i\norm{f_i}_\infty$. Approximate the
continuous coefficients uniformly by trigonometric polynomials.
Each mode $\ee(kt)=\ee(k\beta n)$ combines with the other
phases into an earlier quadratic test. The averages tend to zero
for each fixed trigonometric approximation. Use a countable
approximating family and then let $\delta\downarrow0$.
This proves cancellation for every fixed $h\ne0$ and, by a
countable intersection, simultaneous cancellation for such $h$.

A trigonometric polynomial of $Z_{\rm br}$ has its mean
as the surviving coefficient, and uniform approximation gives the
same comparison for continuous weights. Squeeze $\one_J$
by the same Lipschitz ramps, first controlling the scalar gap.
The prime-count ratio tends to $|J|$ and division gives the
pointwise conclusion.
\end{proof}
\noindent\textbf{Lean status.} The full deduction is checked
under the Heisenberg hypotheses giving Theorem~1.4 on the rotation
extensions and the scalar Vinogradov input. Lean first obtains all
quadratic twists directly from Fourier characters and the smooth
comparison. It then inserts each complex horizontal coefficient by
uniform approximation and bounds the change from the half-open cell
by continuous endpoint ramps. Translation invariance of circle Haar
measure passes the Vinogradov input for $\beta p$ to $\beta(p-1)$.
This proves (56), followed by the continuous and interval selector
comparisons. The composed theorem takes the same literature premises
as the Heisenberg result. Lemma~\ref{checked-prime-distribution}
supplies the scalar Vinogradov estimate.

\section{Separating local and global phase tests}

\begin{proposition}[A.1, separating family]\label{p:separating-family}
\lean{NilBohr.SeparatingFamilyClaims,NilBohr.appendixA_separating_family_of_quantitativeLeibman,NilBohr.separatingFamily_of_obstructions,NilBohr.appendixA_separating_family}\leanok
\uses{selector,phase-tests}
Fix nonconstant Lipschitz $F$ and put $G_q(y)=F(Z(qy))-\mu_F$.
There are sequences $q,M,N=qM^2$ in the smooth-modulus range
with $\log q\asymp(\log r_j)^2$ such that
$\norm{G_q}_{u^3[M]}\gg_F r_j^{-1}$.
For fixed $A,B>0$ and every $P=I\cap(u+q^b\Z)$ with
$I\subset[M]$, $0\le b\le B$ and $|P|\ge(qM)^{1/2}$, the
bound $\norm{G_q}_{u^3[P]}\ll_{A,B,F}(\log M)^{-A}$ holds.
Also, for each fixed integer $S\ge3$, $\norm{G_q}_{U^S[M]}\gg_{F,S}1$.
\end{proposition}
\source{p:separating-family}
\begin{proof}[Proof architecture]
\proves{p:separating-family}
\uses{p:heisenberg-phase,l:heisenberg-horizontal,l:vertical-mode-multiplier,checked-separating-finite,checked-separating-scales,checked-separating-local,checked-separating-quadratic,checked-separating-gowers}\leanok
\paragraph{One sequence of scales.}
Write $r=r_j$ and $R=r_{j+1}$. Take
\[
 \begin{aligned}
 w&=\lfloor\log r\rfloor,&
 q&=\mathcal L(w)^{\left\lfloor w^2/\log\mathcal L(w)\right\rfloor},&
 H&=\left\lfloor R^{1/4}/r\right\rfloor,\\
 M&=rH,& X&=qM,& N&=qM^2.
 \end{aligned}
\]
The formula for $q$ is used for large $j$, where $w\ge16$ and
$\mathcal L(w)\ge2$. The floor gives $w^2-\log\mathcal L(w)<\log q
\le w^2$. Since $\log\mathcal L(w)\le w\log w$, one has
$\log q\sim w^2$. The lower bound on the exponent makes
$8\prod_{p\le w}p\mid q$, and prime divisors of $q$ are $\le w$.
In particular $q$ is coprime to the prime $r>w$.
The growth of $R$ implies that each fixed power of $\log r$ is
$o(\log R)$. As $0\le R^{1/4}-M<r$, it follows that
$M\sim R^{1/4}$. One can use the finite bounds
$H\ge R^{1/6}$, $R^{1/5}\le M\le R^{1/4}$ and
$R^{1/5}\le X\le R^{1/3}$ for large $j$. These give $q\le M$,
$r\le X^{1/100}<R$, and $q^{b+1}\le X^{1/10000}$ uniformly
for $b$ in any fixed finite range. The asymptotic for $M$ also gives $\log\log M\sim\sqrt{\log r}$,
so $\log q\asymp(\log\log M)^4$. Since $N=qM^2$, one has
$w=o((\log N)^c)$ for each fixed $c>0$ and
$\log q=o(\sqrt{\log N})$. The first comparison gives
$w\le(\log\widetilde Q_N)^{c_{\mathrm{rel}}}$ for each fixed
$c_0,c_{\mathrm{rel}}>0$, using $\log\widetilde Q_N\ge(\log N)^{c_0}$.
The second yields $N^{1/2}\exp(-2\sqrt{\log N})\le M\le N^{1/2}$,
since $M=N^{1/2}q^{-1/2}$ exactly. In particular
$M=\lfloor(N/q)^{1/2}\rfloor$.

\paragraph{A correlation on the denominator multiples.}
Fix $h\in\Z\setminus\{0\}$ with $\widehat F(h)\ne0$. Such an $h$
exists because circle characters span a uniformly dense subspace
of the continuous functions. Put $E=\beta-a_j/r$, so
$0<E\le D/R$ with $D$ fixed. For $r\mid y$ in $[M]$,
$a_j qy/r$ is integral and $0\le E qy\le E X<1$ for large $j$.
Consequently $\lfloor\beta qy\rfloor=a_j qy/r$, and
\[
 Q(y)=-\frac{\alpha a_j}{2r}(q^2y^2+qy),\qquad
 |Z(qy)-Q(y)|\le \frac{\abs{\alpha}}2 E X^2.
\]
The last bound tends to zero uniformly because $X\le R^{1/3}$.
Lipschitz continuity therefore permits replacing $F(Z(qy))$
by $F(Q(y))$ in the correlation on $r\mid y$.

For every fixed $k\in\Z\setminus\{0\}$, the phase
$k(Q(r(t+v))-Q(rt))$ is linear in $t$ with coefficient
$-k\alpha a_jq^2rv$. The integer multiplying $\alpha$ is nonzero because
$0<a_j<r$ and $v\ne0$. The geometric identity bounds the
absolute normalized sum over $H$ consecutive integers by
$O(1/(H\tn{k\alpha a_jq^2rv}))$.
Bad approximability of $\alpha$ bounds this by
$O_\alpha(|k|a_jq^2r |v|/H)$.
For each fixed nonzero $v$ this tends to zero, since
$a_j<r$, $H\ge R^{1/6}$, and $q^2r^2=R^{o(1)}$.

For the displayed quantitative bound, replace each unit phase $z_t$ by
$L^{-1}\sum_{a=0}^{L-1}z_{t+a}$, extending the formula for $z_t$
beyond the averaging interval. The average changes by at most
$2L/H$. Apply Cauchy--Schwarz to the smoothed sequence and expand
the square. The diagonal contributes $1/L$ and each off-diagonal
is bounded by $O_\alpha(|k|a_jq^2r L/H)$. Taking
$L=\lfloor H^{1/3}\rfloor$ gives
\[
 \abs{\E_{t\in[H]}\ee(kQ(rt))}^{2}
 \ll_\alpha H^{-1/3}+|k|a_jq^2rH^{-2/3}.
\]
For every $\gamma>0$, $a_jq^2r/H^\gamma\to0$ at the chosen
scales, so this bound tends to zero. The same limit also follows
from the checked argument that takes $j\to\infty$ at fixed $L$
and then lets $L\to\infty$. Weyl's criterion
follows by uniform approximation by trigonometric polynomials.
Apply it to $u\mapsto(F(u)-\mu_F)\ee(-hu)$ to get
\[
 \E_{t\in[H]}G_q(rt)\ee(-hQ(rt))\longrightarrow \widehat F(h).
\]
Inserting $\one_{r\mid y}$ in the average on $[M]$ multiplies this
by $H/M=1/r$. Expanding the indicator into $r$ linear characters
with total coefficient modulus one gives
$\norm{G_q}_{u^3[M]}\gg_F r^{-1}$, since each resulting test is
still quadratic.

\paragraph{Uniform local tests.}
Parametrize a tested $P$ as $a+q^b[|P|]$. The selector then runs
over a progression of step $q^{b+1}$ in the $Z$ variable. The
step conditions established above and the given $|P|\ge X^{1/2}$
allow the centered estimate of Proposition~4.3. The affine change
of variable preserves the degree bound on the test polynomial.
Since $\log X\ge\log M>0$ for large $j$, the resulting bound
$O((\log X)^{-A})$ gives the claimed bound in $\log M$.

\paragraph{Full Gowers norms.}
Choose the fixed Lipschitz cutoff
$\chi(u)=\min(1,\max(0,3\tn u-1))$. Its support lies in
$[1/3,2/3]$, and it is positive near $1/2$, so its circle
integral is positive. Let
$D_h(u)=(F(u)-\mu_F)\ee(-hu)$, with
$\mu_{D_h}=\widehat F(h)$. The chart test with orbit values
$\chi(\{\beta n\})D_h(Z(n))$ splits into the centered
cutoff test for $D_h$ and $\mu_{D_h}$ times the horizontal cutoff.
Both are Lipschitz for the relevant fixed quotient distance.
Use $v=q$ and $K=M\ge X^{1/2}$ in the scalar estimate.
The centered term has zero chart integral, while the horizontal
term has chart mean $\mu_{D_h}\int_\T\chi$. Thus
\[
 \E_{y\in[M]}G_q(y)\chi(\{\beta qy\})\ee(-hZ(qy))
 \longrightarrow \widehat F(h)\int_\T\chi .
\]
The normalized Gowers mean bound and Lemma~4.5 control the
left-hand modulus by
$O_S(\norm{\widehat\chi}_{\ell^1}\norm{G_q}_{U^S[M]})$ for every
fixed $S\ge3$. The limit is nonzero, which gives a positive
lower bound independent of $j$.
\end{proof}
\noindent\textbf{Lean status.} Proposition~A.1 is checked on one
explicit sequence. The local upper bound and full Gowers lower bound use the
specializations of quantitative Leibman. The ordinary phase lower
bound follows from the Diophantine hypotheses, finite Fourier
extraction and the quadratic Weyl argument.

\begin{lemma}[The two dilation bounds and the direct estimate]\label{checked-separating-comparison}
\lean{NilBohr.SeparatingCountComparison,NilBohr.appendixA_count_comparison,NilBohr.separatingFixedFamily_bounds_of_bounded,NilBohr.separating_fixedDilationBound_tendsto,NilBohr.repeatedDilationPhaseBound_one}\leanok
\uses{p:separating-family,t:main,external-inputs,checked-separating-finite}
Use the shifts $y,qy^2$ and $1$-bounded spatial inputs supported
on $[-N,N]$. With $K\ge1$ from the fixed-family theorem, the
dilated upper bound is proportional to
\[
 q^2\left((qM)^{-1}+q^{-1}\norm{G_q}_{u^3[M]}\right)^{1/K},
\]
or to the same expression with $q$ in place of $q^2$ when the
spatial inputs repeat on the $q$ residue classes. The first expression
diverges on the separating scales. The repeated expression diverges
for $K>1$, and tends to zero at $K=1$. The direct estimate of
Theorem~4.2 gives $|T_{N,M,G_q}(\mathbf f)|\ll_F (\log M)^{-1}$
uniformly in the inputs on those scales.
\end{lemma}
\begin{proof}
\proves{checked-separating-comparison}\leanok
Extend the parameter weight and each spatial input by zero off
$q\Z$. For the fixed shifts $n,n^2$, the changes $m=qx$ and $n=qy$
give the same numerator and denominator $(qN)(qM)$. The identity
$\norm{\widetilde G}_{u^3[qM]}=q^{-1}\norm{G_q}_{u^3[M]}$ follows
by transporting quadratic tests in both directions.
Repeating each $f_i(qx+a)=f_i(x)$ for $0\le a<q$ multiplies the
numerator by $q$. The resulting inputs have support in
$[-2qN,2qN]$, since for $|x|\le N$ one has
$|qx+a|\le qN+q\le2qN$ for $N\ge1$. The fixed-family estimate
applies at parameter length $qM$ and spatial scale $(qM)^2=qN$.
Normalizing the bounded weight costs a constant depending on $F$.

For $b=1,2$ write the resulting expression as
$q^b((qM)^{-1}+q^{-1}\norm{G_q}_{u^3[M]})^{1/K}$.
Proposition~A.1 gives a lower bound by a positive constant times
$q^{b-1/K}r^{-1/K}$, which diverges when $b>1/K$ because
$\log r=o(\log q)$. When $b=K=1$, the expression is exactly
$M^{-1}+\norm{G_q}_{u^3[M]}$. The local estimate in
Proposition~A.1 applies to $[M]$ itself since $M\ge X^{1/2}$.

For the direct estimate normalize $G_q$ to a $1$-bounded
weight, and take $\rho=(\log M)^{-1}$ in Theorem~4.2.
The shifts have the required $(1,q)$-coefficients.
A tested interval has length at least
$Mq^{-A_1}(\log M)^{-A_2}$, so its progression has length at least
$Mq^{-A_1-b}(\log M)^{-A_2}-1$. Since $q=M^{o(1)}$ and $b$ is
bounded, this exceeds $X^{1/2}$ for large $j$.
Use Proposition~A.1 with $A\eta\ge1$ and absorb the spatial
term $N^{-\eta}$ into $(\log M)^{-1}$. This gives the uniform bound.
\end{proof}
\noindent\textbf{Lean status.} The comparison is checked from
the inputs of Proposition~A.1, the KMTT inverse estimate and
Ter\'av\'ainen's fixed-family estimate. The divergence statements
concern the displayed upper-bound expressions. The direct
estimate gives the uniform decay of the counts.


\section{Additional claims and scope of the estimates}

\begin{lemma}[1.6, fixed local factors]\label{checked-smooth-factors}
\lean{NilBohr.auxiliary_step_dvd_fixed_power,NilBohr.localGeneralizedVonNeumann_on_lcm_moduli,NilBohr.centeredSelectorComparison_on_lcm_moduli_of_obstructions}\leanok
On the smooth moduli with $\mathcal L(w)\mid q$ reached by the
density iteration, the centered count is $O(w^{-A})$ for a fixed
$A>0$, uniformly over the families and inputs of Lemma~4.4.
\end{lemma}
\begin{proof}[Proof architecture]
\proves{checked-smooth-factors}\uses{external-inputs,checked-difference,p:heisenberg-phase}\leanok
An inverse auxiliary denominator $r\le w$ divides $q$.
Each resulting step $rq^b$ divides the common chosen power
$q^{B+1}$. Use the fixed-step correlations in the spatial
differencing argument and apply the centered phase bound
after substituting the polynomial shifts on the parameter cells.
\end{proof}

\begin{lemma}[Ambient counting comparison]\label{checked-ambient-discrete}
\lean{NilBohr.ambientSelectorCount_eq_explicit,NilBohr.ambientSelectorCount_centered,NilBohr.ambientSelectorCount_sharp,NilBohr.heisenbergAmbientSelectorCount_sharp_of_literature}\leanok
For the shifts $n,n^2$, set $M=\lfloor\sqrt N\rfloor$ and let
$\mathcal C(F)$ be the Mangoldt weighted count with selector test
$F(Z(n))$. For each fixed Lipschitz $F$,
$\mathcal C(F)-\mu_F\mathcal C(1)=o(1)$ uniformly in $A\subset[N]$.
For each fixed open arc $J$ the same comparison holds with
$F=\one_J$ and $\mu_F=|J|$.
The smooth comparison uses Proposition~6.2 and Ter\"av\"ainen's
fixed-family estimate as inputs. The sharp comparison also uses
the prime number theorem.
\end{lemma}
\begin{proof}[Proof architecture]
\proves{checked-ambient-discrete}\uses{p:prime-phase,external-inputs,checked-ergodic}\leanok
Use the prime phase bound and Ter\"av\"ainen's fixed-family
estimate with parameter length $M$ and spatial scale $M^2$.
The change to the ambient normalization costs a factor $M^2/N\le1$.
For interval cutoffs use upper and lower Lipschitz approximants and
the uniform bound on the count with test $1$.
\end{proof}

\begin{lemma}[Ambient Gowers obstructions]\label{checked-ambient-norm}
\lean{NilBohr.ambient_centered_gowers_prescribed_cutoff,NilBohr.tendsto_ambient_centered_localized_mode,NilBohr.ambient_centered_gowers_lower_bound,NilBohr.irrational_quadratic_gowers_two_tendsto_zero,NilBohr.quadratic_obstruction_common_window,NilBohr.no_uniform_weighted_gowers_two}\leanok
For the paper's $\alpha,\beta$ and Lipschitz $F$, prescribe
$h\ne0$ with $\widehat F(h)\ne0$ and a nonzero nonnegative
Lipschitz cutoff $\chi$ supported on an interval of length $<1/2$
inside $(0,1)$. Then, for all sufficiently large $M$,
\[
 0<|\widehat F(h)|\int_{\T}\chi
 \ll \norm{\widehat\chi}_{\ell^1}\norm{F(Z(\cdot))-\mu_F}_{U^3[M]}.
\]
The implicit constant is universal. In particular a nonconstant
$F$ gives nondecay at every fixed order $S\ge3$.
For any irrational $\theta$, the weight $\ee(-\theta n^2)$
has $U^2[M]\to0$, yet has count $1/2$ for suitable $1$-bounded
inputs supported on $[2M^2]$ for the shifts $n,n^2$.
The quadratic example is checked from irrationality alone.
\end{lemma}
\begin{proof}[Proof architecture]
\proves{checked-ambient-norm}\uses{checked-separating-gowers,checked-separating-quadratic}\leanok
At $X=M$ and $q=1$, scalar equidistribution gives the localized
correlation with limit $\widehat F(h)\int\chi$. The seam cutoff
equals one on the prescribed support. Apply the checked localized
multiplier and the order-three bound for the mean. For the quadratic example the
order-two power is a sum of squared linear exponential sums
over interval overlaps. The geometric bound holds away from a
small set of directions. The spatial phases cancel the quadratic
weight exactly on the common supporting count.
\end{proof}

\begin{lemma}[Pointwise limits with Bohr weights]\label{checked-bohr-weights}
\lean{NilBohr.ae_bohrPrimeMode_tendsto,NilBohr.ae_continuousBohrPrimeAverage_tendsto,NilBohr.ae_primeAverage_congruence_tendsto,NilBohr.ae_primeAverage_reduced_congruence_tendsto,NilBohr.ae_simultaneousBohrIntervalAverage_tendsto,NilBohr.simultaneousBohrIntervalDensity_tendsto,NilBohr.ae_selectedSimultaneousBohrIntervalAverage_tendsto}\leanok
For the polynomial ergodic averages, fixed linear twists,
continuous weights of finitely many linear torus coordinates,
fixed congruence classes, and simultaneous interval cutoffs have
almost-everywhere limits at the unrestricted prime normalization.
For reduced congruence classes, and for interval conditions with
positive limiting relative density, the selected-prime normalization
also converges. The proof uses the unrestricted prime pointwise theorem
on the product system, the prime number theorem and Vinogradov.
Reduced congruence normalization uses the progression form of the
prime number theorem.
\end{lemma}
\begin{proof}[Proof architecture]
\proves{checked-bohr-weights}\uses{external-inputs,checked-prime-distribution}
\leanok
Insert a circle rotation in an input at a linear shift, adding that
shift with constant input when needed. Then use the unrestricted
pointwise theorem and Fubini. Finite Fourier sums treat periodic
weights. Uniform trigonometric approximation treats continuous
weights, and boundary strips plus exact rational-grid interpolation
treat intervals, with either endpoint convention.
\end{proof}

\begin{lemma}[The introductory Bohr eigenfunction example]\label{checked-bohr-eigenfunction}
\lean{NilBohr.symmetricBohr_selected_exponential_tendsto,NilBohr.bohr_rotation_eigenfunction_example}\leanok
Let $Tt=t+\theta$ on $\T$ with $\theta$ irrational, $f(t)=\ee(t)$,
$P(n)=n$, and $0<\varepsilon<1/2$. The unrestricted prime average
of $f(T^{p-1}t)$ tends to zero. The average over primes with
$\tn{\theta(p-1)}<\varepsilon$, divided by their count, tends to
$\ee(t)\sin(2\pi\varepsilon)/(2\pi\varepsilon)\ne0$ for every $t$.
\end{lemma}
\begin{proof}
\proves{checked-bohr-eigenfunction}\uses{checked-prime-distribution}\leanok
Write the summand as $\ee(t)\ee(\theta(p-1))$. Vinogradov's
equidistribution gives zero for the unrestricted character average.
For the restricted average, approximate the indicator of the
open symmetric arc from above and below by continuous functions.
The prime density tends to $2\varepsilon$, while the restricted
character sum divided by the unrestricted prime count tends to
$\int_{-\varepsilon}^{\varepsilon}\ee(u)\,du=\sin(2\pi\varepsilon)/\pi$.
The ratio gives the stated limit. It is nonzero because
$0<2\pi\varepsilon<\pi$.
\end{proof}
\noindent\textbf{Lean status.} The open-arc restriction, prime
normalization and two differing limits are checked using the
Vinogradov interface in Lemma~\ref{checked-prime-distribution}.

\begin{lemma}[Progressions in selected primes]\label{checked-prime-progressions}
\lean{NilBohr.fixedWindow_primeProgression_of_sparse,NilBohr.greenTaoRelativeDensity_of_sparse,NilBohr.shrinkingWindow_primeProgression}\leanok
For fixed $J$, positive limiting relative prime density implies a
nontrivial progression of every length $k\ge2$. For $N\ge100$ and
$k\ge4$, the selected primes through
$N$ contain such a progression if their relative density exceeds
$C_k\exp(-(\log\log\log N)^{c_k})$, by Ter\"av\"ainen--Wang
Theorem~1.1. The constants depend only on $k$ and the latter
implication tests the actual density with the arc at scale $N$.
\end{lemma}
\begin{proof}[Proof architecture]
\proves{checked-prime-progressions}\uses{t:pointwise,external-inputs}\leanok
For the first claim, apply the finite density bound as $N\to\infty$,
testing length $\max(k,4)$ and taking a prefix of length $k$.
In the finite claim apply the density bound to the selected prime set.
At $k=4$ its stronger upper bound $(\log\log N)^{-a}$ is at most
$e^{1/(4a)}\exp(-(\log\log\log N)^{1/2})$.
\end{proof}

\begin{lemma}[Restricted maximal domination]\label{checked-maximal-domination}
\lean{NilBohr.selectedPrimeAverage_domination,NilBohr.selectedPrimeMaximal_domination,NilBohr.eLpNorm_selectedPrimeMaximal_le,NilBohr.memLp_selectedErgodicPrimeMaximal,NilBohr.selectedErgodicPrimeMaximal_congr_ae}\leanok
A fixed window with a positive relative-prime density limit has a
constant $C>0$ so that for every summand sequence and every cutoff,
the modulus of its selected-prime average is bounded by $C$
times the unrestricted prime average of the moduli. The same
constant transfers an $L^p$ estimate for the positive unrestricted
maximal operator to the selected maximal operator, including
$p=\infty$.
\end{lemma}
\begin{proof}[Proof architecture]
\proves{checked-maximal-domination}\uses{t:pointwise}\leanok
Compare positive sums and use the ratio of prime counts whenever
the selected set is nonempty. The limiting density bounds the ratio
eventually, and the finitely many earlier cutoffs give a finite
additional constant. An empty selected set has zero average.
Take the supremum in the extended nonnegative reals. Monotonicity
of $L^p$ norms gives the conditional maximal estimate. Measurability
and invariance of the ergodic maxima under a.e. replacement follow
from the countable set of iterates and cutoffs.
\end{proof}


\paragraph{Localized tests versus full Gowers norms.}
The same vertical-mode correlation as in Proposition~A.1 works at
$q=1$. In Lemma~\ref{checked-ambient-norm} the scale $X=M$ is
bracketed by Lemma~6.1, and step $1$ is coprime to every prime
denominator. A fixed short horizontal cutoff with positive Haar
mean and a nonzero Fourier coefficient of $F$ therefore give the
positive lower bound for all sufficiently large $M$.

For the independent order-two example write $I_h=[M]\cap([M]-h)$.
The unnormalized fourth power of the weight is
\[
 \sum_{|h|<M}
 \left|\sum_{x\in I_h}\ee(2\theta hx+\theta h^2)\right|^2 .
\]
Fix a block length $K\ge2$. Irrationality gives a positive least
torus distance to zero among $2\theta j$ for $0<|j|<K$.
For a small enough threshold $\eta>0$, each block of $K$
directions has at most one $h$ with $\tn{2\theta h}<\eta$.
These directions cost at most $(2M/K+1)M^2$ in total. For the
other directions the geometric sum is bounded by $\eta^{-1}$.
The fourth power of the indicator norm is at least $M^3/8$,
as seen by using $0\le h\le M/2$. Divide by that norm,
let $M$ tend to infinity, and then send $K$ to infinity.
For the count, support the base input on $[M^2]$ with
phase $\ee(-\theta x)$. Use the constant and the phase
$\ee(\theta x)$ on $[2M^2]$ as the two shifted inputs.
All parameters $1\le n\le M$ leave the shifted points in their
supports, and the three phase factors multiply to one.
The resulting count normalized at $N=2M^2$ is $1/2$.

\paragraph{Smooth local factors.}
In Remark~1.6 the divisibility $\mathcal L(w)\mid q$ bounds
the cost of prescribing one common step. A localized weighted
correlation is obtained by differencing the spatial inputs and
applying the unweighted KMTT inverse estimate. If the auxiliary
denominator $r$ is at most $w$, then $r\mid q$ and every step
$rq^b$ divides the fixed $q^{B+1}$. On progressions of this
step, approximate each dual factor by its value at the initial
point. Summing the resulting constant correlations gives the
linear-phase extraction at frequency zero. The existing
replacement and deletion argument applies with this fixed-step
localization. The centered progression estimate then bounds the
count uniformly on the iteration's range. In Lean the inverse
threshold is a fixed power of $w^{-A}$, and $A$ is chosen so the
required lcm cutoff is at most $w$.

\paragraph{The ambient sharp-arc comparison.}
For $M=\lfloor\sqrt N\rfloor$ and $A\subset[N]$, the count in (10) is
\[
 \mathcal C(F)=\frac1{NM}\sum_{x\in\Z}\sum_{n=1}^M
  \Lambda(n+1)F(Z(n))\one_A(x)\one_A(x+n)\one_A(x+n^2).
\]
Use Proposition~6.2 at $q=1$
on the full interval $[M]$. A centered Lipschitz selector times
$\Lambda(n+1)$ has arbitrary inverse-logarithmic phase savings.
Divide the weight by a constant times $\log M$ to make it
$1$-bounded. Ter\"av\"ainen's fixed-family estimate uses spatial
scale $M^2$. The inputs on $[N]$ fit in its window of size a
constant times $M^2$, and the base support makes enlarging the
summation window exact. Rescale the count by $M^2/N\le1$ and
restore the weight factor, choosing the original logarithmic
saving large enough to absorb it. Every bound is uniform in
$A$. For an open arc, sandwich its indicator by fixed Lipschitz
functions with means approaching $|J|$. The prime number theorem
bounds $\mathcal C(1)$ independently of $A$ and makes the
sandwich error arbitrarily small.

\paragraph{Pointwise comparison for Bohr weights.}
A single imposed linear twist $\ee(\omega n)$ can be
realized by adjoining a circle rotation by $\omega/a$
and the character in the input at the linear shift $an$.
If necessary add a linear shift with constant input and
reindex the polynomials and functions. Pointwise convergence
on the product system and Fubini then give the twisted limit.
The Lean proof uses the existing two-circle extension with
one rotation frequency set to zero.
For periodic weights expand into finitely many such modes.
For continuous tests on a finite torus, take uniform trigonometric
approximants, intersect their countably many conull sets, and
use their uniform error and bounded inputs.

At each irrational coordinate, Vinogradov bounds the limiting
relative number of primes visiting a small strip around the
two endpoints. At each rational coordinate, Fourier interpolation
on its finite grid reproduces the indicator exactly, including
whatever endpoint values the interval takes. The difference of
the products of the cutoff values is bounded by the sum of the
coordinate errors, so simultaneous conditions require only the
marginal equidistribution statements. Taking the inputs to be
one gives existence of the selected-prime relative density
limit. If it is positive, division by that density changes
from total prime normalization to selected-prime normalization.
For reduced residue classes the prime number theorem in
progressions gives density $1/\varphi(q)$.

\subsection{Fast prime-denominator approximation}

\begin{proposition}[Introductory counterexample]\label{checked-intro-fast-obstruction}
\uses{selector}
\lean{NilBohr.intro_fast_approximation_counterexample}\leanok
There are $\beta\in(0,1)\setminus\mathbb Q$ and strictly increasing
prime denominators $r_j$ with $0<a_j<r_j$ such that
$0<\beta-a_j/r_j<\exp(-r_j^j)$ for every $j$. Set
$N_j=\lfloor\exp(r_j^j/6)\rfloor$. For each $\alpha\in\R$ there are
sets $A_j\subset[N_j]$ with the following properties.
If $x,x+n,x+n^2$ are in $A_j$ with $n\ge1$, then
$Z(n)\bmod 1\notin J=(1/4,1/2)$. The scales tend to infinity and
the density satisfies $|A_j|/N_j\ge 1/(578r_j)$ eventually.
For every fixed $\kappa>0$, the same sets satisfy
$|A_j|/N_j>(\log N_j)^{-\kappa}$ eventually.
\end{proposition}

\begin{proof}
\proves{checked-intro-fast-obstruction}\leanok
Start with $a_0/r_0=1/3$. Choose each prime $r_{j+1}$ with
$r_{j+1}\ge2r_j$ and $r_{j+1}>2\exp(r_j^j)$.
The rational $a_{j+1}/r_{j+1}$ is the first point of that grid
strictly above $a_j/r_j$. Each increment is at most $1/r_{j+1}$.
The geometric tail bound gives a limit $\beta\le2/3$ and
\[
 0<\beta-a_j/r_j\le 2/r_{j+1}<\exp(-r_j^j).
\]
For each fixed $m$, this error is below $r_j^{-m}$ for large $j$,
because $r_j^j\ge m\log r_j$ eventually. Lean uses Mathlib's
Liouville-number criterion for irrationality.

Fix $\alpha$ and consider the $r_j$-multiples in $[N_j]$. A
configuration on this grid has $r_j\mid n$ and $n\le N_j$.
The error $(\beta-a_j/r_j)n$ is less than one, since
$N_j\exp(-r_j^j)\le \exp(-5r_j^j/6)<1$.
It follows that $\lfloor\beta n\rfloor=a_jn/r_j$, so
$Z(n)=c_1n+c_2n^2$ with $c_1=-\alpha\beta/2$ and
$c_2=\alpha(\beta/2-a_j/r_j)$.
Partition the two coordinates $(c_1x,c_2x)\bmod 1$ into a
$17$ by $17$ grid of half-open boxes. A largest cell $A_j$
contains at least $\lfloor N_j/r_j\rfloor/17^2$ points.
For a configuration within the cell, the first coordinate at
$x+n$ and $x$ gives $\tn{c_1n}<1/17$, while the second at
$x+n^2$ and $x$ gives $\tn{c_2n^2}<1/17$.
The torus triangle inequality gives $\tn{Z(n)}<2/17<1/8$.
Every point of $J$ has torus norm greater than $1/4$, proving
avoidance for all positive parameters, including $n=p-1$ for prime $p$.

Eventually $N_j\ge r_j$ and
$\lfloor N_j/r_j\rfloor\ge N_j/(2r_j)$, giving the
claimed density. The floor in $N_j$ yields
$\log N_j\ge r_j^j/12$ eventually.
For fixed $\kappa>0$, take $j$ large enough that $j\kappa\ge2$.
Then $(\log N_j)^{-\kappa}\le12^\kappa r_j^{-2}$,
which is eventually below $1/(578r_j)$ as $r_j\to\infty$.
The sets were chosen before $\kappa$.
\end{proof}

\noindent\textbf{Lean status.} The approximants, exact floor
identity, finite coloring, and eventual density comparison are checked
without external analytic inputs. The same $\beta$ works for every
$\alpha$, including the badly approximable case used in the manuscript.

\paragraph{Progressions, maximal estimates, and available scales.}
For a moving arc $J_N$, the finite progression implication tests
the actual density through $N$. The manuscript gives no uniform
lower bound on this density over moving arcs. In the case $k=4$,
put $L=\log\log\log N\ge0$ and write the stronger density threshold
as $e^{-aL}$. The inequality $\sqrt L-aL\le1/(4a)$ verifies
the weaker common form for the same finite cutoff.
For the maximal deduction, take absolute values of the inputs
before forming the unrestricted multilinear maximal expression.
Lemma~\ref{checked-maximal-domination} works at every cutoff,
so taking the supremum transfers any available estimate for
this positive maximal expression with the same input exponents.

In the density iteration the logarithmic modulus is at most
$b^{i_*}\log\mathcal L(w)$ and the geometric density increase gives
$i_*\le1+\log(1/\delta)/\log(1+\eta)$.
Using $\log\mathcal L(w)\le w\log w\le w^2$ and
$w\le\lceil K\delta^{-E}\rceil$ gives
$\log q_i\ll\delta^{-D_0}$ with $D_0=2E+\log b/\log(1+\eta)$.
If $\delta\ge(\log\log N)^{-c}$ with $cD_0<1$, this bound is
$o(\log\log N)$ uniformly. Since
$\log\log N_i\ge(3/4)\log\log N$, for every $\epsilon>0$ all stages
satisfy $q_i\le(\log N_i)^\epsilon$ when $N$ is large enough.
For a loss $q_i^K$, take $\epsilon<(A-a)/K$ when $K>0$ to pass
from logarithmic saving $A$ to $a<A$. For the other density regime
the witnesses in Lemma~\ref{checked-iteration-moduli} have
$w^2/2\gg(\log N)^{2\kappa E}$ and hence
$\log q>A\log\log N$ for every fixed $A$. These witnesses show
what the admissible range permits at a given scale.

\section{Manuscript and Lean correspondence}
\printmapping

\subsection{Checked Lean steps}

These entries are proofs of the specified steps. They supplement
the manuscript-to-Lean statement map. In particular, an encoded
target proposition with an open proof has no \verb|\leanok| tag.

\begin{lemma}[Finite norms and Lipschitz circle Fourier series]\label{checked-norm-foundations}
\lean{NilBohr.intervalGowersForm_norm_le,NilBohr.normalizedGowers_triangle,NilBohr.normalizedGowersAdjacentMonotonicity,NilBohr.NormalizedGowersLaws.checked,NilBohr.progressionGowers_one_eq_norm_mean,NilBohr.progressionGowers_le_of_le,NilBohr.phaseNorm_le_progressionGowers,NilBohr.normalizedGowers_dilation,NilBohr.kmttNormalizedDilation,NilBohr.bernsteinCircleLipschitzFourier,NilBohr.fourierCoeff_tail_eta_le,NilBohr.norm_fourierSeries_truncation_eta_le}\leanok
The interval cube norm obeys the triangle inequality for every positive
order. On every nonempty finite arithmetic progression, normalized
order comparison and $u^s\ll_s U^s$ have uniform constants. In the
change of variable $x=qy+b$ the normalized dilation bound loses at
most an order-dependent constant times $q^{(s+1)/2^s}$. Lipschitz
circle cutoffs have absolute Fourier mass bounded by a universal
constant times their supremum norm plus Lipschitz seminorm.
\end{lemma}
\begin{proof}
\proves{checked-norm-foundations}\leanok
Coordinate Cauchy--Schwarz bounds the mixed cube form. A separate
direction-sum Cauchy--Schwarz and indicator cube counts give the
normalized order comparison. Exact affine cube transport handles
progressions, and a finite residue Fourier expansion handles dilation.
Parseval applied to translated differences, followed by a limit on
each finite frequency set, bounds the weighted Fourier square sum.
Cauchy--Schwarz against the reciprocal frequencies gives absolute
summability. The reciprocal-square tail is at most $4/R$ outside
$|h|\le R$, giving an $O(\operatorname{Lip}(F)\eta)$ absolute
Fourier tail at the manuscript's cutoff $R=\eta^{-2}$.
\end{proof}

\begin{lemma}[Fixed local factors from the inverse theorem]\label{checked-weak-regularity}
\lean{NilBohr.KMTTLocalInverse.toLocalWeakRegularity,NilBohr.exists_localInverse_fixed_coordinate,NilBohr.fixedStep_progression_partition}\leanok
The local inverse estimate supplies the weak-regularity interface
used in Proposition~3.1, including the partition, progression length
and small-scale alternative.
\end{lemma}
\begin{proof}
\proves{checked-weak-regularity}\uses{external-inputs}\leanok
Choose a fixed power of $q$ divisible by every eligible inverse step.
Partition its residue classes into nearly equal blocks at the
coordinate-dependent resolutions. A telescoping discrepancy is witnessed
by a normalized residual in one coordinate. Reanchor at $P_1$ when
the residual is in the base coordinate. Each atom has zero residual
mean, so a function with the inverse Lipschitz bound has correlation
less than the lower bound forced by the inverse theorem.
\end{proof}

\begin{lemma}[Equivalent prime number theorem limits]\label{checked-prime-normalization}
\lean{NilBohr.PrimeNumberTheorem.of_theta,NilBohr.PrimeNumberTheorem.of_psi,NilBohr.PrimeNumberTheorem.of_count,NilBohr.PrimeNumberTheorem.to_psi,NilBohr.logarithmicIntegral_normalization_tendsto,NilBohr.PrimeNumberTheoremInProgressions.of_logarithmicIntegral}\leanok
The theta, psi and normalized ordinary prime-count asymptotics give
equivalent premises for the conversion to ordinary prime averages.
The fixed-modulus prime count can use either $\operatorname{Li}(N)$
or $N/\log N$ as its scale.
\end{lemma}
\begin{proof}
\proves{checked-prime-normalization}\leanok
The contribution of higher prime powers is $o(N)$. For any fixed
$0<a<1$, the primes at most $N^a$ contribute $o(1)$ to the normalized
count, and on the tail $a\log N\le\log p\le\log N$. Taking $a$ close
to one proves the equivalence with the ordinary count limit.
The same split for the integral of $1/\log t$ on $[2,N]$ proves the
equivalence of the fixed-modulus scales.
\end{proof}

\begin{lemma}[The Heisenberg return window]\label{checked-heisenberg-window}
\lean{NilBohr.heisenbergWindow,NilBohr.isOpen_heisenbergWindow,NilBohr.heisenbergWindow_measure,NilBohr.heisenbergRotation_mem_window_iff,NilBohr.heisenbergWindow_basicNilsequence,NilBohr.heisenbergWindowNilsequence_re_pos_iff,NilBohr.heisenbergWindowNilsequence_prime_re_pos_iff}\leanok
The window $0<y<1$, $z\bmod 1\in J$ defines an open subset of the
right-coset quotient and has Haar mass $|J|$. For admissible $\beta$
and each nonzero integer $n$, the orbit at time $n$ belongs to this
window exactly when $Z(n)\bmod 1\in J$. A basic polynomial nilsequence
realizes the window by the open right half-plane in $\C$.
\end{lemma}
\begin{proof}
\proves{checked-heisenberg-window}\uses{selector}\leanok
Pull the window back to the lattice translates of the open strip.
In the unit Haar fundamental domain its volume is $|J|$ by Fubini.
The signed group powers give the orbit formula. A nonzero integer
multiple of $\beta$ avoids the seam by irrationality. The exact
orbit membership criterion includes the seam condition at time zero.
The distance to the closed complement is a continuous function on
the quotient, and the signed powers are polynomial for the rational
lower central filtration. Complexifying the distance supplies the
required continuous test.
\end{proof}

\begin{lemma}[Prime distribution from the model and sieve]\label{checked-prime-distribution}
\lean{NilBohr.tendsto_cramerWeight_mean,NilBohr.ambient_prime_polynomial_test_error_tendsto,NilBohr.primeNumberTheorem_of_prime_model,NilBohr.primeNumberTheoremInProgressions_of_kmtt_and_sieve,NilBohr.tendsto_cramer_linear_mean,NilBohr.vinogradov_of_prime_model}\leanok
Banks's zero exclusion, the ambient Gowers model, Mertens' interval
bound and the fundamental sieve lemma supply the ordinary prime
number theorem, its fixed reduced progression form, and irrational
linear equidistribution over primes against continuous tests.
\end{lemma}
\begin{proof}
\proves{checked-prime-distribution}\uses{external-inputs,checked-norm-foundations,checked-prime-normalization}\leanok
Average the sieve bounds over a complete period to get the exact
Euler center. On the actual interval or the fixed reduced residue
class each divisor count has error at most one. The order-two
Gowers error transfers the means to Mangoldt, including the fixed
residue projection. For an irrational frequency, use reduced
Dirichlet approximants with denominator at most $\lfloor\sqrt N\rfloor$.
Compare the complex sieve midpoint with its normalized rational
Fourier coefficient over a complete period. Off the multiples of
the denominator the geometric sums are short. On the multiples,
the residual phase differs from its block mean by a controlled
error. The periodic coefficient is bounded by $\sqrt{2/q}$ and the
approximate denominator $q$ tends to infinity at all large scales.
Transfer the linear tests through the model and apply trigonometric
density after conversion to ordinary primes.
\end{proof}

\begin{lemma}[Modulus budgets in Appendix A]\label{checked-iteration-moduli}
\lean{NilBohr.log_iteration_modulus_le_density_power,NilBohr.eventually_density_iteration_moduli_small,NilBohr.eventually_small_modulus_absorbed_by_logarithmic_saving,NilBohr.eventually_large_admissible_modulus_at_polylog_density}\leanok
Given the Section~3 iteration budgets, $\log q_i\ll\delta^{-D_0}$
for a fixed $D_0$. If $\delta\ge(\log\log N)^{-c}$ and $cD_0<1$,
then $q_i\le(\log N_i)^\epsilon$ for every fixed $\epsilon>0$,
uniformly over stages. Every fixed power loss in $q_i$ can be
absorbed by any positive slack in a logarithmic saving. At the
polylogarithmic threshold the admissible numerical parameters allow
$q>(\log N)^A$ for each fixed $A\ge0$.
\end{lemma}
\begin{proof}
\proves{checked-iteration-moduli}\leanok
Bound the stage index using the geometric density increase and
use $\log N_i\ge(\log N)^{3/4}$. For the admissibility example
choose $w=\lceil K(\log N)^{\kappa E}\rceil$ with $\kappa E<c_{\rm rel}$
and take the largest power of $\mathcal L(w)$ with logarithm at
most $w^2$. Its logarithm exceeds $w^2/2$, and the prime support
and divisibility requirements follow from those of $\mathcal L(w)$.
\end{proof}

\begin{lemma}[Rescaling and the vertical chart]\label{checked-rescaling}
\lean{NilBohr.eval_rescalePolynomial,NilBohr.eval_rescalePolynomial_mul,NilBohr.HasCoefficients.rescale,NilBohr.heisenberg_chart_right_lattice}
\leanok
For a polynomial with zero constant coefficient, rescaling has the
prescribed coefficients and commutes with successive affine lifts.
The Heisenberg coordinate change under the right lattice is integral
in the vertical phase and preserves the horizontal fractional part.
\end{lemma}
\begin{proof}
\proves{checked-rescaling}\leanok
The polynomial calculation is exact in $\Z[y]$ and preserves degree
for nonzero scaling. The chart statement follows by the group law
and the formula for adding an integer to a floor.
\end{proof}


\begin{lemma}[Existence of admissible Liouville parameters]\label{checked-beta}
\lean{NilBohr.exists_admissibleBeta_between,NilBohr.AdmissibleBeta.liouville}
\leanok
Every nonempty interval inside $(0,1)$ contains an admissible
$\beta$, and each admissible $\beta$ is a Liouville number.
\end{lemma}
\begin{proof}
\proves{checked-beta}\leanok
Choose prime denominators using Bertrand's postulate at the
prescribed scale and round the numerators upward.
The rational approximants increase, while the approximants
plus the reciprocal of their denominators decrease. Their
supremum satisfies the one-sided approximation bounds.
The resulting error is smaller than every fixed inverse power
of the preceding denominator for sufficiently large indices.
\end{proof}

\begin{lemma}[Localized $U^S$ equality]\label{checked-cube}
\lean{NilBohr.cubeDifference_heisenberg_of_short,NilBohr.gowersCubePower_localized_vertical,NilBohr.normalizedGowers_localized_vertical,NilBohr.normalizedGowers_mul_polynomial_phase}
\leanok
For $S\ge3$, the exact equality (45) holds between the norm of
$f\chi$ with and without the localized vertical exponential,
using the finite zero-extended cube-power model. Multiplication
by an ordinary polynomial phase of degree less than $S$ also leaves
the normalized $U^S$ quantity unchanged.
\end{lemma}
\begin{proof}
\proves{checked-cube}\leanok
The finite cube calculation proves the phase factor cancels in every
nonzero term, and the same cube power and normalizer yield equality.
The multiplier proof of Lemma~4.5 combines this equality with the
Fourier expansion and the normalized Gowers triangle inequality.
For an ordinary polynomial, each additive difference lowers the
degree, so its $S$-fold difference vanishes and the exponential cube
factor is one.
\end{proof}

\begin{lemma}[Horizontal Diophantine and floor identities]\label{checked-floor}
\lean{NilBohr.horizontal_nonzero_alpha_bound,NilBohr.horizontal_zero_alpha_bound,NilBohr.bracketPhase_add_heisenbergPhase,NilBohr.floor_mul_eq_int_of_small_error,NilBohr.horizontal_scale_bounds,NilBohr.heisenberg_horizontal}
\leanok
The finite separation inequalities have a main term
$c_\alpha/(|k|vr^2)$ when $k\ne0$ and $1/r$ when $k=0$,
minus the corresponding approximation error $|l|v|\beta-a/r|$.
The bracket carry identity and small positive approximation
floor identity hold under the displayed arithmetic hypotheses.
\end{lemma}
\begin{proof}
\proves{checked-floor}\leanok
The finite separation uses nearest-integer distance and integer
divisibility. The bracket identity follows by splitting at
$\{\beta n\}=1-\beta$. The floor identity follows from the
unit interval above an integer. A second step checks the scale
uniformity for any approximating pair bracketing $X^{1/100}$.
\end{proof}

\begin{lemma}[Scalar Heisenberg equidistribution]\label{checked-scalar}
\lean{NilBohr.scalar_heisenberg_equidistribution,NilBohr.scalar_heisenberg_log_saving}
\leanok
The scalar estimate (25) follows from the fixed-Heisenberg
horizontal-character obstruction.
\end{lemma}
\begin{proof}
\proves{checked-scalar}\uses{checked-floor,external-inputs}\leanok
The finite-dilation lower bound rules out horizontal characters
of polylogarithmic size, even with prefix length $K\ge X^{1/3}$.
The hypothesis supplies the specialized Green--Tao criterion,
including its chart-mean and metric adaptation. The criterion
gives scalar equidistribution under the stated range.
\end{proof}

\begin{lemma}[Count splitting and density step]\label{checked-density}
\lean{NilBohr.selector_count_split,NilBohr.exists_density_increase_of_count_discrepancy,NilBohr.exists_densest_residue_progression,NilBohr.lift_scaled_configuration,NilBohr.density_iteration_contradiction,NilBohr.recurrence_of_density_increment_at_scale,NilBohr.heisenbergRecurrence_of_density_increment,NilBohr.density_increment_atom_from_count_inputs,NilBohr.density_model_count_lower,NilBohr.regularized_count_discrepancy,NilBohr.card_coincidentShiftParameters_le,NilBohr.norm_selected_count_le_exception_mass,NilBohr.log_initialModulus_eq_mangoldt_sum,NilBohr.exists_log_initialModulus_le_linear,NilBohr.eventually_log_initialModulus_at_density_le,NilBohr.exists_log_iteration_modulus_le_linear_in_w}
\leanok
The selector count splits by linearity. Given equal-mass factor
averages with sufficient count discrepancy, one of them
exceeds the prescribed density ceiling. There is a densest residue
progression, and rescaled configurations preserve the selected
prime. The logarithm of the initial modulus is the Mangoldt sum
through $w$, hence $O(w)$. A finite sequence of admissible
fixed-factor density increments contradicts density at most one.
\end{lemma}
\begin{proof}
\proves{checked-density}\uses{checked-rescaling}\leanok
Finite sums give the two-term weight split and product telescoping.
The mass identity gives the $\ell^1$ bound used in (20).
The affine identities lift every configuration point and the prime.
The prime powers dividing $\mathcal L(w)$ are exactly those at most
$w$. The primorial bound and the higher-prime-power estimate give
Chebyshev's $O(w)$ bound, including the initial finite cutoffs.
For the initial density and for every admissible stage this supplies
the linear-in-$w$ modulus budget used in the manuscript.
Logarithmic counting of the number of stages gives the finite
contradiction. The complete finite recursion takes Proposition~3.1
and a numerical parameter condition as premises. The later
asymptotic argument discharges that numerical condition. The
uniform density increment is also checked from Theorem~3.2
and KMTT weak regularity.
\end{proof}

\begin{lemma}[Local phases on one parameter cell]\label{checked-local-phase}
\lean{NilBohr.product_local_linear_phases,NilBohr.norm_local_phase_product_average_le,NilBohr.norm_partition_delete_le,NilBohr.norm_countingForm_le_local_tests_after_replacement}
\leanok
For fixed spatial $x$, the product of the local linear phases
on a single retained progression in $y$ is a polynomial phase
of degree at most $d$. Its average against the test weight is
bounded by the local polynomial-phase norm.
\end{lemma}
\begin{proof}
\proves{checked-local-phase}\leanok
The exponent is the sum of scalar multiples of the shifts and
constant terms. Polynomial degree bounds and finite products
give the stated reduction after the spatial blocks and residues
have been fixed. The finite selection and deletion bounds partition
$[M]$ into intervals of length $H$ and a tail, then count spatial
pairs within the polynomial displacement of a block boundary.
The retained count is bounded by the supremum of the local tests.
\end{proof}

\begin{lemma}[Sieve midpoint on a progression]\label{checked-sieve}
\lean{NilBohr.sieveMidpoint_bound,NilBohr.exists_modEq_of_dvd_affine,NilBohr.abs_card_dvd_affine_sub_le_one,NilBohr.primeProgression_as_affine,NilBohr.primeProgression_sieve_scales,NilBohr.cramerWeight_eq_sieveIndicator,NilBohr.cramerSieve_correlation_bound,NilBohr.primeDivisors_heisenberg_log_saving,NilBohr.eventually_cramerSieve_coefficients,NilBohr.MertensIntervalProductBound.eulerProduct_le_log}
\leanok
With upper and lower divisor envelopes, replacing the indicator by
their midpoint incurs at most half their average gap. For a
divisor coprime to the progression step, divisibility restricts to a
single residue class of parameter indices, whose count differs from
length divided by the divisor by at most one. The change to
$m=qy+1$ preserves the finite average and the degree of its
polynomial test. The resulting sieve progressions have the
uniform length and step bounds stated in Proposition~6.2. On the
residue-one progression the Cram\'er weight is its Euler product
normalization times the coprimality indicator omitting primes that
divide $q$. From the quoted sieve and interval-product assumptions,
the envelopes have absolute reciprocal-weighted gap at most
$(\log X)^{-D}$ for any fixed $D>0$ and all sufficiently large $X$.
\end{lemma}
\begin{proof}
\proves{checked-sieve}\uses{checked-twisted,external-inputs}\leanok
The midpoint inequality is pointwise followed by the triangle
inequality in a finite average. The arithmetic constraint is a
linear congruence with invertible coefficient.
The progression size estimates use the floor in $M$ and the
polylogarithmic bounds on $q$ and the first step. The finite
correlation bound sums the exceptional-divisor contribution using a
harmonic sum. On every other divisor progression that is coprime
to the first step, the twisted estimate applies to vertically
mean-zero tests. Deleting factors preserves the Mertens interval
bound. At level $\lfloor X^{1/20000}\rfloor$, the logarithmic
level ratio $U$ grows at least as a positive power of $\log X$.
The fundamental lemma then gives the stated envelope error.
\end{proof}

\begin{lemma}[Pointwise convergence and normalization steps]\label{checked-ergodic}
\lean{NilBohr.ae_mangoldtAverage_tendsto_zero_of_log_decay,NilBohr.ae_primeAverage_tendsto_zero_of_mangoldt_log_decay,NilBohr.ae_unrestrictedPrimeAverage_tendsto,NilBohr.OpenArc.exists_lipschitz_sandwich,NilBohr.centered_primeAverage_of_sandwich,NilBohr.ae_centered_selector_of_smooth_limits,NilBohr.pointwise_of_smooth_limits,NilBohr.pointwise_of_smooth_log_decay}
\leanok
For logarithmically bounded summands, a prescribed $L^1$ rate
$O((\log H)^{-6})$ implies almost-everywhere convergence to zero
of the prefix averages. Given the scalar counting and logarithmic
mass asymptotics, bounded log-weighted averages and unweighted
averages have the same limit. Scalar gap bounds turn centered lower
approximants into a centered interval weight. A positive prime-density
ratio and a centered comparison identify the selected prime limit.
\end{lemma}
\begin{proof}
\proves{checked-ergodic}\leanok
The first proof uses the manuscript's
$\lfloor\exp((j+1)^{1/3})\rfloor$ sampling and a gap bound.
Prime-power removal uses an explicit count and Mangoldt mass bound.
Log-weight removal uses the mass deficit of the weights $\log p/\log H$
conditional on the two scalar PNT asymptotics. The unrestricted
prime limit then follows conditionally on the weighted
Mirek--Wan--Wright statement. For interval weights the checked
Lipschitz ramps come from distance to the arc centre on the circle,
including arcs that cross the coordinate seam.
The last two deductions squeeze the centered interval average and
divide by the positive limiting count ratio. The full deduction
in Theorem~1.4 from the smooth logarithmic estimate is therefore
checked. The analytic bound at each positive exponent is checked from
Proposition~6.2 and the fixed-family estimate, and the bracket
deduction uses the rotation and cutoff lemmas below.
\end{proof}

\begin{lemma}[Orbit transference and duality]\label{checked-orbit}
\lean{NilBohr.orbit_counting_eq_average,NilBohr.integral_orbit_countingForm,NilBohr.integral_norm_weightedErgodicAverage_le_of_orbit_bound,NilBohr.orbit_bound_of_teravainen,NilBohr.exists_polynomial_orbit_window,NilBohr.mangoldt_L1_log_decay_of_primeLocalPhase}
\leanok
For integrable weighted averages and a common spatial window
containing all shifted orbit arguments, a uniform bound on
the discrete shifted count transfers to the $L^1$ norm of
the weighted ergodic average.
\end{lemma}
\begin{proof}
\proves{checked-orbit}\leanok
Average the measure-preserving change of variables over each
spatial index. Test against $\overline{A(x)}/|A(x)|$ on the
nonzero set of the weighted average $A$ and use zero otherwise.
This isolates the $L^1$ norm. For the fixed-family application the finite count estimate is the
explicit Ter\'av\'ainen Theorem~4.1 input. The common spatial window
uses a sum of absolute polynomial coefficients to bound all shifts.
The rate is supplied by Proposition~6.2 at $q=1$.
\end{proof}

\begin{lemma}[Spatial differencing and exceptional weights]\label{checked-difference}
\lean{NilBohr.countingForm_sq_le_spatial_energy,NilBohr.spatial_slice_square_eq_differences,NilBohr.exclude_exceptional_model_zero,NilBohr.weighted_local_inverse_of_kmtt_degree_uniform}
\leanok
Cauchy--Schwarz removes a bounded parameter weight, and grouping
spatial pairs by their difference gives the exact squared slice
in Lemma~5.1. The separate real-zero threshold comparison uses
the explicit Banks premise and the conductor bound.
\end{lemma}
\begin{proof}
\proves{checked-difference}\uses{external-inputs}\leanok
Use finite complex Cauchy--Schwarz and collect pairs by their
integer difference. For the real-zero step, multiply
$\log r\leq\log\widetilde Q$ by $1-\beta_0>0$.
The ambient prime model is an explicit KMTT input. The degree-uniform
KMTT local inverse hypothesis supplies a common bound before $C$ and
the family length are fixed. Spatial differencing and subsequent
phase extraction preserve that bound.
\end{proof}

\begin{lemma}[Product rotations and the bracket identity]\label{checked-rotations}
\lean{NilBohr.exists_quadratic_completion,NilBohr.ergodicSequence_circleRotationExtension,NilBohr.ae_quadratic_twisted_interval_of_smooth,NilBohr.ae_tendsto_primeAverage_fourier_quadratic_twist,NilBohr.phase_bracket_eq_piecewise,NilBohr.ae_bracket_fourier_cancellation,NilBohr.pointwise_bracket_of_heisenberg_smooth}
\leanok
The enlarged family and the circle-product system recover the
quadratic twist in (54). Equations (53) and (55) follow from
the smooth comparison on the product system. The bracket
carry decomposition preceding (56) also holds exactly.
\end{lemma}
\begin{proof}
\proves{checked-rotations}\uses{checked-ergodic,checked-floor}\leanok
Adjoin the absent degrees and extend the functions by constants.
The product system is probability preserving. Its character
is the desired quadratic twist times a circle factor of modulus
one. Fubini removes that factor. The floor formula gives both
branches of the bracket exponential. Fourier approximation and
endpoint ramps transfer these pieces to prime averages using
Vinogradov's scalar input. This gives (56) and the bracket
pointwise comparison.
\end{proof}


\begin{lemma}[Product obstruction and twisted prefixes]\label{checked-twisted}
\lean{NilBohr.choose_modEq_of_modEq_factorial_mul,NilBohr.realBinomialPhase_exp_totalVariation_le,NilBohr.heisenbergCoverTest_chartMean_eq_zero_of_lipschitz,NilBohr.norm_heisenbergTwistedAverage_le_of_product_obstruction,NilBohr.twisted_heisenberg_log_saving_of_dilated,NilBohr.twisted_heisenberg_log_saving_for_denominators,NilBohr.twisted_heisenberg_phaseNorm}
\leanok
Given the fixed scalar and product character criteria and explicit
bounds on the finite prefix scales, the polynomially twisted
Heisenberg average is bounded in terms of the two obstruction
thresholds and the variation of the residual twist. Under the
dilated scalar and horizontal estimates this yields a uniform
logarithmic saving at every exponent.
\end{lemma}
\begin{proof}
\proves{checked-twisted}\uses{checked-floor,checked-scalar,external-inputs}\leanok
The product criterion isolates a nonzero circle frequency.
Rounding in the integer binomial basis and splitting modulo
$|m|d!$ yields the covering-map expression and a slowly varying
error factor. Haar integration proves zero mean on the vertical
fibres of the covering test. The scalar criterion for each
sufficiently long prefix and partial summation give the finite bound.
The logarithmic scales then absorb the cover's distortion and the
discarded short prefixes. This deduction is for vertically mean-zero
tests on the quotient.
\end{proof}


\begin{lemma}[Selector cutoffs and discrepancy]\label{checked-selector-cutoffs}
\lean{NilBohr.quantitativeSelectorCutoffs,NilBohr.horizontalCircleEstimate_of_scalar,NilBohr.norm_countingForm_localized_selector_le,NilBohr.norm_countingForm_selector_boundary_le,NilBohr.norm_countingForm_selector_discrepancy_le}
\leanok
Three nonnegative short-arc cutoffs and a boundary cutoff give a
partition of unity. The Lipschitz bounds are $O(\varepsilon^{-1})$
and the boundary mean is $O(\varepsilon)$. The finite count
discrepancy splits into the localized prime-versus-model counts
and the boundary count.
\end{lemma}
\begin{proof}
\proves{checked-selector-cutoffs}\uses{checked-scalar,checked-cube,external-inputs}\leanok
Normalize the tent functions by their sum with a positive floor.
Outside the endpoint strip the floor has no effect, and the
Lipschitz and measure bounds follow from the lower bound on the
denominator and the length of a ball. Horizontal projection of the
scalar orbit estimate bounds the mean on the finite progression.
The localized multiplier and normalized Gowers laws then control
the prime discrepancy, using the Bernstein Fourier bound.
\end{proof}


\begin{lemma}[Centered chart test and seam removal]\label{checked-centered}
\lean{NilBohr.heisenbergCenteredTest_orbit,NilBohr.heisenbergCenteredTest_verticalMean_zero,NilBohr.heisenbergCenteredTest_lipschitz,NilBohr.heisenberg_seamAverage_of_horizontal,NilBohr.centeredHeisenbergPhaseEstimate_of_obstructions}
\leanok
A Lipschitz circle weight centered at its Haar mean yields a
right lattice-invariant, vertically mean-zero test after cutting off
the horizontal seam. Its quotient Lipschitz norm is
$O_F(1+\varepsilon^{-1})$. The scalar and twisted orbit estimates
then give Equation~(27) uniformly.
\end{lemma}
\begin{proof}
\proves{checked-centered}\uses{checked-floor,checked-scalar,checked-twisted,external-inputs}\leanok
The orbit formula uses the vertical chart coordinate
$z-x\lfloor y\rfloor$. A relative coordinate bound for each group
increment controls the change in the centered test. If the horizontal
coordinate crosses the chart seam, the cutoff bounds the surviving
values by the horizontal displacement divided by $\varepsilon$.
Summing over chains and taking the quotient infimum gives the
Lipschitz estimate. Circle translation gives the vertical mean, and a
horizontal majorant bounds the discarded seam average.
\end{proof}

\begin{lemma}[Prime-weighted seam error]\label{checked-prime-seam}
\lean{NilBohr.primeProgression_mangoldt_seam_error}\leanok
At the seam width $\eta=(\log X)^{-A-2}$, replacing the centered
selector on the shifted-prime progression by the full cutoff chart
test costs $O_F((\log X)^{-A})$ against $\Lambda$, uniformly in the
polynomial phase.
\end{lemma}
\begin{proof}
\proves{checked-prime-seam}\uses{checked-centered,checked-sieve}\leanok
At divisor one the prime progression has length at least $X^{1/2}$
and step at most $X^{1/10000}$. The common horizontal denominator
is coprime to its step. The scalar estimate bounds the seam mass
by $O(\eta)$. The von Mangoldt weight costs at most $\log X$, and the
polynomial phase has modulus one.
\end{proof}

\begin{lemma}[Cram\'er cancellation with the full chart test]\label{checked-prime-cramer}
\lean{NilBohr.norm_cramerPhaseAverage_le_of_sieve,NilBohr.eventually_siftedPhase_error,NilBohr.primeCramerPhase_log_saving}\leanok
The full centered chart test at seam width $(\log X)^{-E}$,
for any fixed $E>0$, has arbitrary logarithmic saving on the
prime progression against the Cram\'er model, uniformly in the
polynomial test of degree at most $d$.
\end{lemma}
\begin{proof}
\proves{checked-prime-cramer}\uses{checked-sieve,checked-centered}\leanok
Use the exact affine reindexing and normalize the chart test by a
fixed sup bound. On every good divisor progression, all the chart's
vertical fibres have mean zero and its quotient Lipschitz norm
is polylogarithmic. Harmonic summation, the exceptional divisors,
and the sieve gap give one more inverse logarithm than the target,
before multiplication by the Euler product.
\end{proof}

\begin{lemma}[Fourier restriction on a progression]\label{checked-fourier-restriction}
\lean{NilBohr.localized_prime_model_replacement_progression,NilBohr.residue_indicator_fourier_l1,NilBohr.intervalCyclicFourierCoefficient_l1_le,NilBohr.quantitativePrimeSelectorCutoffs,NilBohr.localized_prime_model_replacement_full_chart,NilBohr.prime_model_replacement_centered_chart,NilBohr.primeModelReplacement_log_saving_of_kmtt}
\leanok
A localized and polynomially twisted progression average of an
ambient weight is controlled by its ambient $U^S$ quantity, with the
length ratio $X/|P'|$, the Fourier mass of the circle cutoff, and
a factor bounded by $2+\log X$. The bound extends to the full
centered seam test and yields arbitrary logarithmic saving for the
ambient prime-model discrepancy on the progressions of Proposition~6.2.
\end{lemma}
\begin{proof}
\proves{checked-fourier-restriction}\uses{checked-cube,checked-selector-cutoffs,external-inputs}\leanok
Use cyclic Fourier orthogonality to write the two interval and
congruence indicators as linear phases. The residue coefficients
have mass one. Geometric summation bounds each nonzero interval
coefficient by the reciprocal of twice its centered frequency,
so the total mass is at most one plus the harmonic sum to $X$.
The Gowers triangle and mean bounds and the checked phase
equalities give the normalized restriction.
The bound is uniform in the vertical Fourier index. Absolute Fourier
convergence and an exact short-arc partition of the seam cutoff give a
polylogarithmic loss for the full chart test. The prime-model bound
absorbs this loss.
\end{proof}

\begin{lemma}[Scalar and product Heisenberg models]\label{checked-group-model}
\uses{selector}
\lean{NilBohr.heisenbergLeibmanBasis_rational,NilBohr.heisenbergProductLeibmanBasis_rational,NilBohr.heisenbergNilmanifoldMetric_topology,NilBohr.heisenbergProductNilmanifoldMetric_topology,NilBohr.malcevQuotientChainDistance_eq_heisenbergQuotientDistance,NilBohr.malcevQuotientChainDistance_eq_heisenbergProductQuotientDistance,NilBohr.isLeibmanHaar_heisenbergNilmanifoldMeasure,NilBohr.isLeibmanHaar_heisenbergProductNilmanifoldMeasure,NilBohr.heisenbergNilmanifoldMeasure_unique,NilBohr.heisenbergProductNilmanifoldMeasure_unique,NilBohr.heisenbergTwist_productLipschitzBound}\leanok
The scalar and product quotients have adapted Mal'cev bases of
rationality height at most one. Their chain metrics induce the
quotient topologies. Unit-cube volume pushes forward to the normalized
Haar probabilities and gives the chart integral on each. The
dilated scalar orbit is filtered polynomial of degree at most two.
Adjoining an ordinary polynomial of degree at most $d$ gives the
product of filtration degree $\max(2,d)$. The twist
$\Phi(g)\ee(-u)$ is lattice invariant, has mean zero after circle
integration, and increases the scalar supremum-plus-Lipschitz
bound by at most an absolute factor.
\end{lemma}
\begin{proof}
\proves{checked-group-model}\uses{checked-floor}\leanok
The global matrix charts are Euclidean, with polynomial group laws.
The lattices are closed and discrete and the closed unit cubes
map onto the compact Hausdorff quotients. Smooth one-parameter axes
give second-kind coordinates $(x,y,z-xy)$ and the extra real coordinate.
The sole nonzero scalar Lie bracket of basis axes is
$[X,Y]=Z$ up to sign from exchanging the axes. Ordering the extra axis
before $Z$ when $d=1$ and after $Z$ for $d\ge2$ makes the basis
adapted to the appropriate product filtration.

Small chain displacements control the coordinate differences and
conversely. The quotient infimum over the lattice gives the actual
quotient metric because the lattice is closed. For Haar, take the
pushforward of normalized volume on the half-open fundamental cube.
Shears preserve ambient Lebesgue measure, and periodicity gives the
chart formula. Uniqueness follows by lifting an invariant quotient
probability with counting measure on the lattice fibres and comparing
the resulting ambient Haar measures. The scalar and circle
projections are bounded by the product chain distance, giving the
product test bound by expanding the difference of two products.
\end{proof}

\begin{lemma}[Quantitative Leibman specialization]\label{checked-green-tao-adapter}
\lean{NilBohr.scalarHeisenbergObstruction_of_greenTao,NilBohr.productHeisenbergObstruction_of_greenTao,NilBohr.SelectorLiteratureInputs.of_quantitativeLeibman,NilBohr.PrimePhaseLiteratureInputs.of_quantitativeLeibman}\leanok
The general quantitative Leibman input gives the scalar and product
character criteria used in Proposition~4.3 on the respective
nilmanifolds of Lemma~\ref{checked-group-model}.
The recurrence and prime-phase input bundles can therefore be
assembled directly from the listed source interfaces.
\end{lemma}
\begin{proof}
\proves{checked-green-tao-adapter}\uses{external-inputs,checked-group-model}\leanok
Use the descended scalar test on the dilated orbit and the test
$\Phi(g)\ee(-u)$ on the product orbit. Their means and Lipschitz bounds
are those of the concrete model. In the product application a fixed
factor in the test bound is absorbed by lowering the tolerance.
A horizontal character vanishes on the commutator axis. Its
remaining integer coordinates give the horizontal linear form
and, in the product case, the coefficient of the real polynomial.
First differences in the scalar case and the iterated differences
at zero in the product case give the required binomial-coefficient
inequalities. Constants are chosen after the dimension and
filtration degree and before the orbit and test parameters.
\end{proof}

\begin{lemma}[Finite identities for Appendix A]\label{checked-separating-finite}
\uses{selector,d:Cq}
\lean{NilBohr.phaseNorm_dilatedWeight,NilBohr.parameterLength_separatingFamily,NilBohr.countingForm_separatingFamily_dilate,NilBohr.countingForm_separatingFamily_repeat,NilBohr.residueCorrelation_le_phaseNorm,NilBohr.abs_heisenbergPhase_sub_separatingQuadratic_le}\leanok
Zero extension to $q\Z$ gives a phase norm of exactly $q^{-1}$
times the original phase norm. The fixed-family count normalization
contributes $q^2$, or $q$ when the spatial inputs repeat on all
$q$ residue classes. For a rational approximation with
$0\le\beta-a/r\le E$ and $E qM<1$, the error in the quadratic
approximation to $Z(qy)$ on $r\mid y$ is at most
$|\alpha|E(qM)^2/2$.
\end{lemma}
\begin{proof}
\proves{checked-separating-finite}\uses{checked-floor,checked-fourier-restriction}\leanok
Affine changes of variable transport the real polynomial tests in
both directions. The residue expansion has Fourier mass one, so a
correlation on the multiples of $r$ contributes its value divided
by $r$ to the global phase-norm lower bound. On those multiples
the small nonnegative approximation error fixes the floor exactly.
The count identities follow by reindexing both summation variables
and tracking the numerator and denominator separately.
\end{proof}

\begin{lemma}[The separating scales]\label{checked-separating-scales}
\uses{selector,d:Cq}
\lean{NilBohr.SeparatingScale,NilBohr.eventually_separatingScale,NilBohr.tendsto_separatingLengthQuotient,NilBohr.tendsto_separatingLength,NilBohr.eventually_separatingStep,NilBohr.eventually_separatingRelativeModulus,NilBohr.eventually_separatingRelativeScaleWindow,NilBohr.tendsto_separatingCoefficient_div_lengthQuotient,NilBohr.tendsto_separatingApproximation,NilBohr.eventually_separatingLogModulus_logLogLength,NilBohr.tendsto_logDenominator_div_logSeparatingModulus,NilBohr.tendsto_separatingLogLength_div_exp_sqrt_log,NilBohr.tendsto_logSeparatingScaleX_div_logNext,NilBohr.eventually_initialModulus_dvd_separatingModulus,NilBohr.tendsto_separatingW_div_logSpatialScale_pow,NilBohr.tendsto_logSeparatingModulus_div_sqrt_logSpatialScale}\leanok
The floored power of $\mathcal L(w)$ and the integer multiple
of the denominator in the definition of $M$ give positive
scales, $\log q\asymp(\log r_j)^2$, and
$q\le M$. The logarithmic ratios give
$\log M\sim\tfrac14\exp(\sqrt{\log r_j})$ and
$qM=r_{j+1}^{1/4+o(1)}$. At $X=qM$, the denominator bracket and the bounded
powers of the progression step are valid for large $j$.
\end{lemma}
\begin{proof}
\proves{checked-separating-scales}\uses{checked-separating-finite}\leanok
The floor changes $\log q$ by at most $w\log w$. The next
denominator grows faster than any fixed power of the current
one or of the chosen modulus. The resulting polynomial bounds
in $r_{j+1}$ control $H$, $M$ and $X$ on the same sequence.
\end{proof}

\begin{lemma}[Local tests at the separating scales]\label{checked-separating-local}
\uses{selector,d:Cq}
\lean{NilBohr.SeparatingLocalPhaseDecay,NilBohr.separatingLocalPhase_log_bound,NilBohr.eventually_separating_test_length,NilBohr.separatingCount_log_bound,NilBohr.separatingGlobalPhase_tendsto_zero,NilBohr.separatingLocalPhase_of_obstructions,NilBohr.separatingCount_of_kmtt}\leanok
The centered progression estimate gives arbitrary logarithmic
saving in the local phase tests at the separating scales.
Every progression required by Theorem~4.2 at $\rho=(\log M)^{-1}$
meets the local length threshold. Testing $[M]$ in particular
gives convergence to zero for its ordinary global $u^3$ norm.
\end{lemma}
\begin{proof}
\proves{checked-separating-local}\uses{p:heisenberg-phase,t:main,checked-separating-scales}\leanok
Reparametrize the integer progressions and transport the
polynomial tests along the affine change. The exact cardinality
bound for the residue class loses at most one. The small powers
of $q$ leave the required length margin.
\end{proof}

\begin{lemma}[Continuous circle tests]\label{checked-separating-fourier}
\lean{NilBohr.exists_nonzero_fourierCoeff_of_nonconstant,NilBohr.circleMean_centered_mul_fourier,NilBohr.tendsto_circle_average_of_fourier}\leanok
A nonconstant continuous circle function has a nonzero Fourier
coefficient at some nonzero frequency. The mean of its centered
function times the matching conjugate character is that coefficient.
Cancellation of every nonzero mode on finite rows gives
equidistribution under continuous tests.
\end{lemma}
\begin{proof}
\proves{checked-separating-fourier}\leanok
Use uniqueness of circle Fourier coefficients for the first
claim and trigonometric polynomial approximation for the last.
\end{proof}

\begin{lemma}[Quadratic correlation on the denominator multiples]\label{checked-separating-quadratic}
\uses{selector,d:Cq}
\lean{NilBohr.prefixAverage_tendsto_zero_of_shift_correlations,NilBohr.separatingQuadratic_modes_tendsto_zero,NilBohr.separatingQuadratic_mode_sq_le_growing_shift,NilBohr.tendsto_separatingCoefficient_div_lengthQuotient_pow,NilBohr.separatingGlobalPhase_lower_bound}\leanok
At the separating scales the phase on the denominator multiples
is approximated by a quadratic. Its nonzero Fourier modes have
vanishing averages. The correlation with a quadratic test and
the residue expansion give a positive multiple of $r_j^{-1}$
as a lower bound for the full-interval phase norm.
\end{lemma}
\begin{proof}
\proves{checked-separating-quadratic}\uses{checked-separating-finite,checked-separating-scales,checked-separating-fourier}\leanok
The geometric-series and bad-approximation bounds control each
nonzero shift in the quadratic sequence. Growing-shift van der
Corput gives the displayed $H^{-1/3}+|k|a_jq^2r_jH^{-2/3}$ bound,
whose terms vanish at the separating scales. A second checked
proof takes fixed shifts first. Trigonometric approximation gives
the limit, and the surviving constant is a nonzero Fourier coefficient.
\end{proof}

\begin{lemma}[A full Gowers lower bound]\label{checked-separating-gowers}
\uses{selector,d:Cq}
\lean{NilBohr.separating_gowers_lower_bound}\leanok
For every fixed $S\ge3$, scalar Heisenberg equidistribution on
the long, coprime progression and the localized vertical
multiplier give a positive lower bound for $U^S[M]$.
\end{lemma}
\begin{proof}
\proves{checked-separating-gowers}\uses{p:heisenberg-phase,l:vertical-mode-multiplier,checked-separating-scales,checked-separating-fourier}\leanok
The product of the horizontal cutoff mean and a nonzero
vertical Fourier coefficient is the limit of the tested
average. The Gowers mean bound and Lemma~4.5 bound the
average by a constant times the full Gowers norm.
\end{proof}

